\nonstopmode{}
\documentclass[a4paper]{book}
\usepackage[times,inconsolata,hyper]{Rd}
\usepackage{makeidx}
\makeatletter\@ifl@t@r\fmtversion{2018/04/01}{}{\usepackage[utf8]{inputenc}}\makeatother
% \usepackage{graphicx} % @USE GRAPHICX@
\makeindex{}
\begin{document}
\chapter*{}
\begin{center}
{\textbf{\huge Package `GLLAMMR'}}
\par\bigskip{\large \today}
\end{center}
\ifthenelse{\boolean{Rd@use@hyper}}{\hypersetup{pdftitle = {GLLAMMR: Generalized Linear Latent and Mixed Models}}}{}
\ifthenelse{\boolean{Rd@use@hyper}}{\hypersetup{pdfauthor = {Josh McGrane}}}{}
\begin{description}
\raggedright{}
\item[Type]\AsIs{Package}
\item[Title]\AsIs{Generalized Linear Latent and Mixed Models}
\item[Version]\AsIs{1.2.0}
\item[Description]\AsIs{Comprehensive implementation of Generalized Linear Latent and Mixed
Models following the 'Stata' 'GLLAMM' framework by Rabe-Hesketh, Skrondal,
and Pickles (2004) <}\Rhref{https://doi.org/10.1007/BF02295939}{doi:10.1007/BF02295939}\AsIs{>. Supports multilevel generalized
linear models with Gaussian, binomial, Poisson, ordinal, and multinomial
responses; item response theory models including dichotomous (Rasch, 2PL,
3PL) and polytomous (GRM, PCM, GPCM, NRM) models; differential item
functioning analysis; explanatory item response models with item covariates;
latent class analysis; structural equation models; mixed response models;
and survival models. All models support frequency and probability weights
for survey data and aggregated observations. Provides marginal
(population-averaged) predictions via Monte Carlo integration. Uses
'TMB' (Template Model Builder) for efficient computation via automatic
differentiation and Laplace approximation. Provides comprehensive
diagnostics, visualization, and model comparison tools.}
\item[License]\AsIs{GPL-3}
\item[Depends]\AsIs{R (>= 4.0.0)}
\item[Imports]\AsIs{TMB (>= 1.9.0), Matrix (>= 1.5.0), stats, methods, graphics,
grDevices, utils}
\item[LinkingTo]\AsIs{TMB, RcppEigen}
\item[Suggests]\AsIs{testthat (>= 3.0.0), knitr, rmarkdown, ggplot2, lme4, mirt,
TAM, MASS, lavaan, npmlreg, glmmTMB, ordinal, poLCA, survival,
HSAUR3, ltm}
\item[Encoding]\AsIs{UTF-8}
\item[RoxygenNote]\AsIs{7.3.3}
\item[VignetteBuilder]\AsIs{knitr}
\item[NeedsCompilation]\AsIs{yes}
\item[Author]\AsIs{Josh McGrane [aut, cre]}
\item[Maintainer]\AsIs{Josh McGrane }\email{drjoshmcgrane@gmail.com}\AsIs{}
\end{description}
\Rdcontents{Contents}
\HeaderA{GLLAMMR-package}{GLLAMMR: Generalized Linear Latent and Mixed Models}{GLLAMMR.Rdash.package}
\aliasA{GLLAMMR}{GLLAMMR-package}{GLLAMMR}
\keyword{internal}{GLLAMMR-package}
%
\begin{Description}
Comprehensive implementation of Generalized Linear Latent and Mixed Models following the 'Stata' 'GLLAMM' framework by Rabe-Hesketh, Skrondal, and Pickles (2004) \Rhref{https://doi.org/10.1007/BF02295939}{doi:10.1007\slash{}BF02295939}. Supports multilevel generalized linear models with Gaussian, binomial, Poisson, ordinal, and multinomial responses; item response theory models including dichotomous (Rasch, 2PL, 3PL) and polytomous (GRM, PCM, GPCM, NRM) models; differential item functioning analysis; explanatory item response models with item covariates; latent class analysis; structural equation models; mixed response models; and survival models. All models support frequency and probability weights for survey data and aggregated observations. Provides marginal (population-averaged) predictions via Monte Carlo integration. Uses 'TMB' (Template Model Builder) for efficient computation via automatic differentiation and Laplace approximation. Provides comprehensive diagnostics, visualization, and model comparison tools.
\end{Description}
%
\begin{Author}
\strong{Maintainer}: Josh McGrane \email{drjoshmcgrane@gmail.com}

\end{Author}
\HeaderA{.val\_row}{Build one validation result row}{.val.Rul.row}
\keyword{internal}{.val\_row}
%
\begin{Description}
Build one validation result row
\end{Description}
%
\begin{Usage}
\begin{verbatim}
.val_row(
  case,
  statistic,
  gllammr,
  reference,
  tolerance,
  relative = TRUE,
  note = ""
)
\end{verbatim}
\end{Usage}
\HeaderA{abilities}{Extract Person Abilities from IRT Models}{abilities}
%
\begin{Description}
Extract estimated person abilities (theta) from IRT models.
For multi-level models, this returns the person-level deviations (theta\_0).
Use \code{composite\_theta} to get total abilities including random effects.
\end{Description}
%
\begin{Usage}
\begin{verbatim}
abilities(object, ...)
\end{verbatim}
\end{Usage}
%
\begin{Arguments}
\begin{ldescription}
\item[\code{object}] A fitted IRT model
(theta\_0 + random effects) instead of just theta\_0. Default FALSE.

\item[\code{...}] Additional arguments (not used)
\end{ldescription}
\end{Arguments}
%
\begin{Value}
A named vector of person abilities
\end{Value}
%
\begin{Examples}
\begin{ExampleCode}
## Not run: 
# Person-level deviations
abilities(fit)

# Total abilities (including class effects)
abilities(fit, composite = TRUE)

## End(Not run)

\end{ExampleCode}
\end{Examples}
\HeaderA{abilities.gllamm\_irt\_multilevel}{Total abilities for multi-level IRT fits}{abilities.gllamm.Rul.irt.Rul.multilevel}
%
\begin{Description}
Total abilities for multi-level IRT fits
\end{Description}
%
\begin{Usage}
\begin{verbatim}
## S3 method for class 'gllamm_irt_multilevel'
abilities(object, composite = FALSE, ...)
\end{verbatim}
\end{Usage}
%
\begin{Arguments}
\begin{ldescription}
\item[\code{object}] A fitted multi-level IRT model

\item[\code{composite}] If TRUE, return person deviation plus group effects
(total ability); if FALSE (default), the person deviation only

\item[\code{...}] Additional arguments (not used)
\end{ldescription}
\end{Arguments}
%
\begin{Value}
A named vector of person abilities
\end{Value}
\HeaderA{aghq}{Adaptive quadrature integration specification}{aghq}
%
\begin{Description}
Use as \code{gllamm(..., integration = aghq(k))} to integrate the random
intercept by adaptive Gauss-Hermite quadrature with \code{k} nodes
instead of the Laplace approximation. Currently supports two-level
random-intercept models with gaussian, binomial, or poisson families.
Laplace (the default) is equivalent to \code{aghq(1)}.
\end{Description}
%
\begin{Usage}
\begin{verbatim}
aghq(k = 15)
\end{verbatim}
\end{Usage}
%
\begin{Arguments}
\begin{ldescription}
\item[\code{k}] Number of quadrature nodes (default 15)
\end{ldescription}
\end{Arguments}
%
\begin{Value}
An object of class \code{gllamm\_integration}
\end{Value}
\HeaderA{binomial}{Binomial Family for Binary and Binomial Outcomes}{binomial}
%
\begin{Description}
Create a family object for binomial regression models with various link functions
\end{Description}
%
\begin{Usage}
\begin{verbatim}
binomial(link = c("logit", "probit", "cloglog"))
\end{verbatim}
\end{Usage}
%
\begin{Arguments}
\begin{ldescription}
\item[\code{link}] Link function to use:
\begin{description}

\item["logit"] Logistic regression (default) - symmetric S-curve
\item["probit"] Probit regression - based on normal distribution
\item["cloglog"] Complementary log-log - asymmetric, suitable for rare events and survival data

\end{description}

\end{ldescription}
\end{Arguments}
%
\begin{Details}
The binomial family is used for binary (0/1) or binomial (count/total) responses.

\strong{Logit Link (default):}
\deqn{P(Y=1|x) = \frac{1}{1 + \exp(-x'\beta)}}{}
This is the standard logistic regression. The logit link is symmetric and
appropriate when the probability of success and failure are equally likely
to change with covariates.

\strong{Probit Link:}
\deqn{P(Y=1|x) = \Phi(x'\beta)}{}
where \eqn{\Phi}{} is the standard normal CDF. This link is also symmetric
and yields similar results to logit but with slightly different tail behavior.

\strong{Complementary Log-Log (cloglog) Link:}
\deqn{P(Y=1|x) = 1 - \exp(-\exp(x'\beta))}{}
This link is \emph{asymmetric} and is particularly useful for:
\begin{itemize}

\item{} Rare events (when P(Y=1) is small)
\item{} Survival analysis with discrete time intervals
\item{} Gompertz or extreme value distributions
\item{} When the hazard is proportional (as in survival models)

\end{itemize}


The cloglog link arises naturally when modeling grouped survival data
or when events follow a Poisson process over time intervals.
\end{Details}
%
\begin{Value}
A family object of class \code{binomial\_family} with components:
\begin{ldescription}
\item[\code{family}] Character: "binomial"
\item[\code{link}] Character: name of link function
\item[\code{link\_code}] Integer: numeric code for TMB (1-3)
\end{ldescription}
\end{Value}
%
\begin{SeeAlso}
\code{\LinkA{gllamm}{gllamm}}, \code{\LinkA{ordinal}{ordinal}}
\end{SeeAlso}
%
\begin{Examples}
\begin{ExampleCode}
## Not run: 
# Logistic regression (default)
family1 <- binomial()
family2 <- binomial(link = "logit")

# Probit regression
family3 <- binomial(link = "probit")

# Complementary log-log for rare events
family4 <- binomial(link = "cloglog")

# Use with gllamm() - recommended interface
fit <- gllamm(outcome ~ age + treatment + (1 | clinic),
              data = mydata,
              family = binomial(link = "logit"))

# Rare event with cloglog
fit_rare <- gllamm(rare_disease ~ exposure + (1 | region),
                   data = epi_data,
                   family = binomial(link = "cloglog"))

## End(Not run)

\end{ExampleCode}
\end{Examples}
\HeaderA{coef.gllamm\_eirt}{Extract Item Parameters from EIRT Model}{coef.gllamm.Rul.eirt}
%
\begin{Description}
Extract item parameters with standard errors (if available)
\end{Description}
%
\begin{Usage}
\begin{verbatim}
## S3 method for class 'gllamm_eirt'
coef(object, se = FALSE, ...)
\end{verbatim}
\end{Usage}
%
\begin{Arguments}
\begin{ldescription}
\item[\code{object}] A gllamm\_eirt object

\item[\code{se}] Logical; include standard errors?

\item[\code{...}] Additional arguments (ignored)
\end{ldescription}
\end{Arguments}
%
\begin{Value}
A data frame with item parameters
\end{Value}
\HeaderA{coef.gllamm\_irt}{Extract Coefficients from IRT Models}{coef.gllamm.Rul.irt}
%
\begin{Description}
Extract item parameters from fitted IRT models
\end{Description}
%
\begin{Usage}
\begin{verbatim}
## S3 method for class 'gllamm_irt'
coef(object, type = c("item", "person"), ...)
\end{verbatim}
\end{Usage}
%
\begin{Arguments}
\begin{ldescription}
\item[\code{object}] A fitted IRT model

\item[\code{type}] Type of coefficients: "item" for item parameters,
"person" for person abilities, "random" for random effects (multi-level only)

\item[\code{...}] Additional arguments passed to specific methods
\end{ldescription}
\end{Arguments}
%
\begin{Value}
Item parameters (data frame), person abilities (vector),
or random effects (list)
\end{Value}
%
\begin{Examples}
\begin{ExampleCode}
## Not run: 
# Item parameters
coef(fit, type = "item")

# Person abilities
coef(fit, type = "person")

# Random effects (multi-level only)
coef(fit, type = "random")

## End(Not run)

\end{ExampleCode}
\end{Examples}
\HeaderA{compare\_eirt}{Compare Explanatory IRT Models}{compare.Rul.eirt}
%
\begin{Description}
Compare two or more EIRT models using likelihood ratio tests and information criteria
\end{Description}
%
\begin{Usage}
\begin{verbatim}
compare_eirt(..., test = c("LRT", "none"))
\end{verbatim}
\end{Usage}
%
\begin{Arguments}
\begin{ldescription}
\item[\code{...}] Two or more gllamm\_eirt objects to compare

\item[\code{test}] Type of test: "LRT" for likelihood ratio test, "none" for just IC comparison
\end{ldescription}
\end{Arguments}
%
\begin{Details}
This function compares EIRT models to test the importance of item covariates.
Models are compared using:
\begin{itemize}

\item{} Log-likelihood
\item{} AIC and BIC
\item{} Likelihood ratio test (if nested models)

\end{itemize}


For nested models (e.g., with and without a predictor), the LRT tests
whether adding the predictor significantly improves fit.
\end{Details}
%
\begin{Value}
A data frame with model comparison statistics
\end{Value}
%
\begin{Examples}
\begin{ExampleCode}
## Not run: 
# Fit model with predictor
fit1 <- fit_eirt(responses, item_data,
                 difficulty_formula = ~ word_freq,
                 model = "Rasch")

# Fit model without predictor
fit0 <- fit_eirt(responses, item_data,
                 difficulty_formula = ~ 1,
                 model = "Rasch")

# Compare models
compare_eirt(fit0, fit1)

## End(Not run)

\end{ExampleCode}
\end{Examples}
\HeaderA{compute\_dif\_effect\_size}{Compute effect size for DIF}{compute.Rul.dif.Rul.effect.Rul.size}
\keyword{internal}{compute\_dif\_effect\_size}
%
\begin{Description}
Compute effect size for DIF
\end{Description}
%
\begin{Usage}
\begin{verbatim}
compute_dif_effect_size(fit1, fit2, item)
\end{verbatim}
\end{Usage}
%
\begin{Arguments}
\begin{ldescription}
\item[\code{fit1}] Fitted model for group 1

\item[\code{fit2}] Fitted model for group 2

\item[\code{item}] Item index
\end{ldescription}
\end{Arguments}
%
\begin{Value}
Effect size (ETS Delta scale)
\end{Value}
\HeaderA{compute\_irt\_reliability}{Compute IRT reliability}{compute.Rul.irt.Rul.reliability}
\keyword{internal}{compute\_irt\_reliability}
%
\begin{Description}
Compute IRT reliability
\end{Description}
%
\begin{Usage}
\begin{verbatim}
compute_irt_reliability(object)
\end{verbatim}
\end{Usage}
\HeaderA{compute\_item\_fit\_sx2}{Compute S-X\textasciicircum{}2 item fit statistic}{compute.Rul.item.Rul.fit.Rul.sx2}
\keyword{internal}{compute\_item\_fit\_sx2}
%
\begin{Description}
Compute S-X\textasciicircum{}2 item fit statistic
\end{Description}
%
\begin{Usage}
\begin{verbatim}
compute_item_fit_sx2(object)
\end{verbatim}
\end{Usage}
\HeaderA{compute\_multinomial\_probs}{Compute multinomial probabilities}{compute.Rul.multinomial.Rul.probs}
\keyword{internal}{compute\_multinomial\_probs}
%
\begin{Description}
Compute multinomial probabilities
\end{Description}
%
\begin{Usage}
\begin{verbatim}
compute_multinomial_probs(X, beta, eta_random, n_categories)
\end{verbatim}
\end{Usage}
%
\begin{Arguments}
\begin{ldescription}
\item[\code{X}] Fixed effects design matrix

\item[\code{beta}] Beta matrix (K-1) × p

\item[\code{eta\_random}] Random effects contribution (vector of length n\_obs)

\item[\code{n\_categories}] Number of categories
\end{ldescription}
\end{Arguments}
%
\begin{Value}
Matrix of probabilities (n\_obs × n\_categories)
\end{Value}
\HeaderA{compute\_person\_fit\_outfit\_infit}{Compute outfit/infit person fit statistics}{compute.Rul.person.Rul.fit.Rul.outfit.Rul.infit}
\keyword{internal}{compute\_person\_fit\_outfit\_infit}
%
\begin{Description}
Compute outfit/infit person fit statistics
\end{Description}
%
\begin{Usage}
\begin{verbatim}
compute_person_fit_outfit_infit(object)
\end{verbatim}
\end{Usage}
\HeaderA{compute\_test\_information\_summary}{Compute test information summary}{compute.Rul.test.Rul.information.Rul.summary}
\keyword{internal}{compute\_test\_information\_summary}
%
\begin{Description}
Compute test information summary
\end{Description}
%
\begin{Usage}
\begin{verbatim}
compute_test_information_summary(object)
\end{verbatim}
\end{Usage}
\HeaderA{construct\_Z\_matrix}{Construct random effects design matrix for newdata}{construct.Rul.Z.Rul.matrix}
\keyword{internal}{construct\_Z\_matrix}
%
\begin{Description}
Construct random effects design matrix for newdata
\end{Description}
%
\begin{Usage}
\begin{verbatim}
construct_Z_matrix(newdata, random_terms, group_var = NULL)
\end{verbatim}
\end{Usage}
%
\begin{Arguments}
\begin{ldescription}
\item[\code{newdata}] New data frame

\item[\code{random\_terms}] Parsed random effects terms from original model

\item[\code{group\_var}] Name of grouping variable
\end{ldescription}
\end{Arguments}
%
\begin{Value}
Random effects design matrix Z
\end{Value}
\HeaderA{cooks.distance.gllamm}{Group-level Cook's distance for GLLAMM models}{cooks.distance.gllamm}
%
\begin{Description}
Influence of each cluster on the fixed effects, computed by refitting the
model with the cluster deleted: D\_j = (beta - beta\_(-j))' V\textasciicircum{}(-1)
(beta - beta\_(-j)) / p, with V the estimated covariance of the fixed
effects. Case deletion at the cluster level is the standard influence
measure for mixed models; observation-level deletion would break the
random-effects structure.
\end{Description}
%
\begin{Usage}
\begin{verbatim}
## S3 method for class 'gllamm'
cooks.distance(model, max_groups = 50, ...)
\end{verbatim}
\end{Usage}
%
\begin{Arguments}
\begin{ldescription}
\item[\code{model}] A fitted \code{gllamm} object (from \code{gllamm()})

\item[\code{max\_groups}] Refuse to run for more clusters than this (each cluster
costs one model refit); raise the limit explicitly for large data

\item[\code{...}] Additional arguments (currently unused)
\end{ldescription}
\end{Arguments}
%
\begin{Value}
Named vector of Cook's distances, one per cluster. Clusters whose
deletion refit fails get \code{NA}.
\end{Value}
\HeaderA{create\_grouping\_matrix}{Create grouping factor matrix}{create.Rul.grouping.Rul.matrix}
\keyword{internal}{create\_grouping\_matrix}
%
\begin{Description}
Create grouping factor matrix
\end{Description}
%
\begin{Usage}
\begin{verbatim}
create_grouping_matrix(terms, data)
\end{verbatim}
\end{Usage}
%
\begin{Arguments}
\begin{ldescription}
\item[\code{terms}] Parsed random effects terms

\item[\code{data}] Person-level data
\end{ldescription}
\end{Arguments}
%
\begin{Value}
List with group\_ids matrix and group information
\end{Value}
\HeaderA{diagnostics}{Diagnostic Methods for GLLAMM Models}{diagnostics}
%
\begin{Description}
Comprehensive diagnostics for model checking
\end{Description}
\HeaderA{dif\_plot}{Plot Item Characteristic Curves by group}{dif.Rul.plot}
%
\begin{Description}
Plot Item Characteristic Curves by group
\end{Description}
%
\begin{Usage}
\begin{verbatim}
dif_plot(dif_result, item, ...)
\end{verbatim}
\end{Usage}
%
\begin{Arguments}
\begin{ldescription}
\item[\code{dif\_result}] Object of class dif\_analysis

\item[\code{item}] Item index to plot

\item[\code{...}] Additional plotting parameters
\end{ldescription}
\end{Arguments}
\HeaderA{dif\_test}{Test for Differential Item Functioning (DIF)}{dif.Rul.test}
%
\begin{Description}
Detect whether item parameters differ systematically across groups
\end{Description}
%
\begin{Usage}
\begin{verbatim}
dif_test(
  irt_fit,
  group,
  items = NULL,
  type = c("both", "uniform", "nonuniform"),
  method = c("lr", "wald", "both"),
  alpha = 0.05
)
\end{verbatim}
\end{Usage}
%
\begin{Arguments}
\begin{ldescription}
\item[\code{irt\_fit}] A fitted IRT model from fit\_irt()

\item[\code{group}] Vector of group indicators (must have exactly 2 unique values)

\item[\code{items}] Vector of item indices to test. If NULL, tests all items.

\item[\code{type}] Type of DIF to test: "uniform" (difficulty), "nonuniform" (discrimination), or "both"

\item[\code{method}] Test method: "lr" (likelihood ratio), "wald", or "both"

\item[\code{alpha}] Significance level for flagging items (default 0.05)
\end{ldescription}
\end{Arguments}
%
\begin{Value}
An object of class \code{dif\_analysis} with components:
\begin{ldescription}
\item[\code{dif\_results}] Data frame with test statistics for each item
\item[\code{flagged\_items}] Vector of item indices that show significant DIF
\item[\code{group\_fits}] List of fitted models for each group
\item[\code{baseline\_fit}] The original fitted model
\end{ldescription}
\end{Value}
%
\begin{Examples}
\begin{ExampleCode}
## Not run: 
# Fit IRT model
fit <- fit_irt(responses, model = "2PL")

# Test for DIF across gender
dif_result <- dif_test(fit, group = gender)
print(dif_result)

# Plot ICCs for flagged items
dif_plot(dif_result, item = 5)

## End(Not run)

\end{ExampleCode}
\end{Examples}
\HeaderA{dif\_test\_with\_data}{Test for DIF with explicit response data}{dif.Rul.test.Rul.with.Rul.data}
%
\begin{Description}
Test for DIF with explicit response data
\end{Description}
%
\begin{Usage}
\begin{verbatim}
dif_test_with_data(
  response_matrix,
  group,
  model = c("Rasch", "2PL", "3PL", "GRM", "PCM", "GPCM"),
  items = NULL,
  type = c("both", "uniform", "nonuniform"),
  method = c("lr", "wald", "both"),
  alpha = 0.05
)
\end{verbatim}
\end{Usage}
%
\begin{Arguments}
\begin{ldescription}
\item[\code{response\_matrix}] Matrix of item responses (persons x items)

\item[\code{group}] Vector of group indicators

\item[\code{model}] IRT model type

\item[\code{items}] Items to test (default: all)

\item[\code{type}] Type of DIF

\item[\code{method}] Test method

\item[\code{alpha}] Significance level
\end{ldescription}
\end{Arguments}
%
\begin{Value}
Object of class \code{dif\_analysis}
\end{Value}
%
\begin{Examples}
\begin{ExampleCode}
## Not run: 
dif_result <- dif_test_with_data(responses, group = gender, model = "2PL")

## End(Not run)

\end{ExampleCode}
\end{Examples}
\HeaderA{drop\_intercept\_column}{Drop the intercept column from a fixed-effects design matrix}{drop.Rul.intercept.Rul.column}
\keyword{internal}{drop\_intercept\_column}
%
\begin{Description}
Cumulative-link (ordinal) models absorb the location into the thresholds;
keeping a free intercept alongside free thresholds leaves the model
unidentified (only their difference enters the likelihood).
\end{Description}
%
\begin{Usage}
\begin{verbatim}
drop_intercept_column(X)
\end{verbatim}
\end{Usage}
%
\begin{Arguments}
\begin{ldescription}
\item[\code{X}] Design matrix from model.matrix()
\end{ldescription}
\end{Arguments}
%
\begin{Value}
X without its "(Intercept)" column
\end{Value}
\HeaderA{eirt\_item\_thresholds}{Reconstruct per-item threshold parameters from a fitted polytomous EIRT model}{eirt.Rul.item.Rul.thresholds}
\keyword{internal}{eirt\_item\_thresholds}
%
\begin{Description}
Mirrors the threshold construction in the TMB template (gllamm\_eirt.hpp):
GRM uses the first-offset + log-spacing parameterization, PCM/GPCM use
sum-to-zero step deviations around the item location, LPCM uses the
threshold regression plus step residuals.
\end{Description}
%
\begin{Usage}
\begin{verbatim}
eirt_item_thresholds(object)
\end{verbatim}
\end{Usage}
%
\begin{Arguments}
\begin{ldescription}
\item[\code{object}] Fitted gllamm\_eirt model (polytomous)
\end{ldescription}
\end{Arguments}
%
\begin{Value}
List of per-item threshold vectors on the absolute scale
\end{Value}
\HeaderA{eirt\_poly\_model\_name}{Map an EIRT poly\_model\_type code to the shared probability helper's model name}{eirt.Rul.poly.Rul.model.Rul.name}
\keyword{internal}{eirt\_poly\_model\_name}
%
\begin{Description}
Map an EIRT poly\_model\_type code to the shared probability helper's model name
\end{Description}
%
\begin{Usage}
\begin{verbatim}
eirt_poly_model_name(poly_model_type)
\end{verbatim}
\end{Usage}
\HeaderA{eirt\_r\_squared}{Compute R-squared for Item Parameter Regression}{eirt.Rul.r.Rul.squared}
%
\begin{Description}
Calculate proportion of variance in item parameters explained by covariates
\end{Description}
%
\begin{Usage}
\begin{verbatim}
eirt_r_squared(object, parameter = c("difficulty", "discrimination"))
\end{verbatim}
\end{Usage}
%
\begin{Arguments}
\begin{ldescription}
\item[\code{object}] A gllamm\_eirt object

\item[\code{parameter}] Which parameter: "difficulty" or "discrimination"
\end{ldescription}
\end{Arguments}
%
\begin{Value}
R-squared value
\end{Value}
\HeaderA{expand\_nested\_random\_term}{Expand a nested random-effects term into one term per level}{expand.Rul.nested.Rul.random.Rul.term}
\keyword{internal}{expand\_nested\_random\_term}
%
\begin{Description}
(x | a/b/c) becomes the lme4-equivalent (x | a) + (x | a:b) + (x | a:b:c).
Each expanded term carries an interaction grouping over the levels above it.
\end{Description}
%
\begin{Usage}
\begin{verbatim}
expand_nested_random_term(parsed_term)
\end{verbatim}
\end{Usage}
%
\begin{Arguments}
\begin{ldescription}
\item[\code{parsed\_term}] Output of parse\_random\_term()
\end{ldescription}
\end{Arguments}
%
\begin{Value}
List of parsed terms (length 1 for non-nested input)
\end{Value}
\HeaderA{expand\_nested\_terms}{Expand nested terms}{expand.Rul.nested.Rul.terms}
\keyword{internal}{expand\_nested\_terms}
%
\begin{Description}
Expand nested terms
\end{Description}
%
\begin{Usage}
\begin{verbatim}
expand_nested_terms(terms, data)
\end{verbatim}
\end{Usage}
\HeaderA{extract\_random\_terms}{Extract random effects terms from formula}{extract.Rul.random.Rul.terms}
\keyword{internal}{extract\_random\_terms}
%
\begin{Description}
Extract random effects terms from formula
\end{Description}
%
\begin{Usage}
\begin{verbatim}
extract_random_terms(formula)
\end{verbatim}
\end{Usage}
\HeaderA{extract\_random\_vcov}{Extract random effects variance-covariance matrix from fitted model}{extract.Rul.random.Rul.vcov}
\keyword{internal}{extract\_random\_vcov}
%
\begin{Description}
Extract random effects variance-covariance matrix from fitted model
\end{Description}
%
\begin{Usage}
\begin{verbatim}
extract_random_vcov(object)
\end{verbatim}
\end{Usage}
%
\begin{Arguments}
\begin{ldescription}
\item[\code{object}] Fitted gllamm object
\end{ldescription}
\end{Arguments}
%
\begin{Value}
Variance-covariance matrix of random effects
\end{Value}
\HeaderA{families}{Family Objects for GLLAMM Models}{families}
%
\begin{Description}
Family constructors for generalized linear latent and mixed models
\end{Description}
\HeaderA{find\_outliers}{Outlier detection for GLLAMM}{find.Rul.outliers}
%
\begin{Description}
Outlier detection for GLLAMM
\end{Description}
%
\begin{Usage}
\begin{verbatim}
find_outliers(object, threshold = 3, ...)
\end{verbatim}
\end{Usage}
%
\begin{Arguments}
\begin{ldescription}
\item[\code{object}] A gllamm object

\item[\code{threshold}] Threshold for standardized residuals (default: 3)

\item[\code{...}] Additional arguments
\end{ldescription}
\end{Arguments}
%
\begin{Value}
Indices of potential outliers
\end{Value}
\HeaderA{fit}{Generic Fit Statistics Function}{fit}
%
\begin{Description}
Compute model-specific fit statistics for GLLAMM objects
\end{Description}
%
\begin{Usage}
\begin{verbatim}
fit(object, ...)
\end{verbatim}
\end{Usage}
%
\begin{Arguments}
\begin{ldescription}
\item[\code{object}] A fitted model object

\item[\code{...}] Additional arguments passed to methods
\end{ldescription}
\end{Arguments}
%
\begin{Details}
The \code{fit()} function provides comprehensive, model-specific fit statistics:

\strong{GLMM models:}
\begin{itemize}

\item{} Log-likelihood, AIC, BIC
\item{} R-squared (marginal and conditional for Gaussian models)
\item{} Intraclass correlation coefficient (ICC)

\end{itemize}


\strong{IRT models:}
\begin{itemize}

\item{} Log-likelihood, AIC, BIC
\item{} Item fit statistics (S-X\textasciicircum{}2)
\item{} Person fit statistics (outfit/infit)
\item{} Reliability estimates
\item{} Test information function

\end{itemize}


\strong{Latent Class Analysis:}
\begin{itemize}

\item{} Log-likelihood, AIC, BIC
\item{} Entropy (classification quality)
\item{} Class proportions
\item{} Average posterior probabilities (APPA)

\end{itemize}


\strong{Ordinal models:}
\begin{itemize}

\item{} Log-likelihood, AIC, BIC
\item{} Pseudo-R\textasciicircum{}2 (McFadden)
\item{} Proportional odds test (for PO/probit models)

\end{itemize}

\end{Details}
%
\begin{Value}
An object of class \code{fit\_statistics} with model-specific components
\end{Value}
%
\begin{Examples}
\begin{ExampleCode}
## Not run: 
# GLMM
fit1 <- gllamm(y ~ x + (1 | group), data = data)
fit(fit1)

# IRT
fit2 <- fit_irt(responses, model = "2PL")
fit(fit2, compute_item_fit = TRUE)

# LCA
fit3 <- fit_lca(indicators, nclass = 3)
fit(fit3)

# Ordinal
fit4 <- fit_ordinal(rating ~ x + (1 | id), data = data)
fit(fit4, test_po = TRUE)

## End(Not run)

\end{ExampleCode}
\end{Examples}
\HeaderA{fit.gllamm}{Fit Statistics for Standard GLLAMM Models}{fit.gllamm}
%
\begin{Description}
Fit Statistics for Standard GLLAMM Models
\end{Description}
%
\begin{Usage}
\begin{verbatim}
## S3 method for class 'gllamm'
fit(object, quiet = FALSE, ...)
\end{verbatim}
\end{Usage}
%
\begin{Arguments}
\begin{ldescription}
\item[\code{object}] A gllamm object

\item[\code{quiet}] Suppress ICC computation messages (default: FALSE)

\item[\code{...}] Additional arguments (not currently used)
\end{ldescription}
\end{Arguments}
%
\begin{Value}
Object of class \code{fit\_gllamm} with fit statistics
\end{Value}
\HeaderA{fit.gllamm\_eirt}{Fit Statistics for Explanatory IRT Models}{fit.gllamm.Rul.eirt}
%
\begin{Description}
Fit Statistics for Explanatory IRT Models
\end{Description}
%
\begin{Usage}
\begin{verbatim}
## S3 method for class 'gllamm_eirt'
fit(object, ...)
\end{verbatim}
\end{Usage}
%
\begin{Arguments}
\begin{ldescription}
\item[\code{object}] A gllamm\_eirt object

\item[\code{...}] Additional arguments
\end{ldescription}
\end{Arguments}
%
\begin{Value}
Object of class \code{fit\_eirt} with EIRT-specific fit statistics
\end{Value}
\HeaderA{fit.gllamm\_irt}{Fit Statistics for IRT Models}{fit.gllamm.Rul.irt}
%
\begin{Description}
Fit Statistics for IRT Models
\end{Description}
%
\begin{Usage}
\begin{verbatim}
## S3 method for class 'gllamm_irt'
fit(object, compute_item_fit = TRUE, compute_person_fit = TRUE, ...)
\end{verbatim}
\end{Usage}
%
\begin{Arguments}
\begin{ldescription}
\item[\code{object}] A gllamm\_irt object

\item[\code{compute\_item\_fit}] Compute item fit statistics (S-X\textasciicircum{}2) (default: TRUE)

\item[\code{compute\_person\_fit}] Compute person fit statistics (outfit/infit) (default: TRUE)

\item[\code{...}] Additional arguments
\end{ldescription}
\end{Arguments}
%
\begin{Value}
Object of class \code{fit\_irt} with IRT-specific fit statistics
\end{Value}
\HeaderA{fit.gllamm\_lca}{Fit Statistics for Latent Class Analysis}{fit.gllamm.Rul.lca}
%
\begin{Description}
Fit Statistics for Latent Class Analysis
\end{Description}
%
\begin{Usage}
\begin{verbatim}
## S3 method for class 'gllamm_lca'
fit(object, ...)
\end{verbatim}
\end{Usage}
%
\begin{Arguments}
\begin{ldescription}
\item[\code{object}] A gllamm\_lca object

\item[\code{...}] Additional arguments
\end{ldescription}
\end{Arguments}
%
\begin{Value}
Object of class \code{fit\_lca} with LCA-specific fit statistics
\end{Value}
\HeaderA{fit.gllamm\_multinomial}{Fit Statistics for Multinomial Regression Models}{fit.gllamm.Rul.multinomial}
%
\begin{Description}
Fit Statistics for Multinomial Regression Models
\end{Description}
%
\begin{Usage}
\begin{verbatim}
## S3 method for class 'gllamm_multinomial'
fit(object, ...)
\end{verbatim}
\end{Usage}
%
\begin{Arguments}
\begin{ldescription}
\item[\code{object}] A gllamm\_multinomial object

\item[\code{...}] Additional arguments
\end{ldescription}
\end{Arguments}
%
\begin{Value}
Object of class \code{fit\_multinomial} with multinomial-specific fit statistics
\end{Value}
\HeaderA{fit.gllamm\_ordinal}{Fit Statistics for Ordinal Regression Models}{fit.gllamm.Rul.ordinal}
%
\begin{Description}
Fit Statistics for Ordinal Regression Models
\end{Description}
%
\begin{Usage}
\begin{verbatim}
## S3 method for class 'gllamm_ordinal'
fit(object, test_po = TRUE, ...)
\end{verbatim}
\end{Usage}
%
\begin{Arguments}
\begin{ldescription}
\item[\code{object}] A gllamm\_ordinal object

\item[\code{test\_po}] Test proportional odds assumption (default: TRUE for logit/probit)

\item[\code{...}] Additional arguments
\end{ldescription}
\end{Arguments}
%
\begin{Value}
Object of class \code{fit\_ordinal} with ordinal-specific fit statistics
\end{Value}
\HeaderA{fit\_binomial}{Fit Binomial Regression Models with Random Effects}{fit.Rul.binomial}
%
\begin{Description}
Fit logistic, probit, or complementary log-log models for binary or binomial
responses. This function can be called directly or through \code{gllamm()} with
\code{family = binomial()}. The \code{gllamm()} interface is recommended
for consistency with other model types.
\end{Description}
%
\begin{Usage}
\begin{verbatim}
fit_binomial(
  formula,
  data,
  link = c("logit", "probit", "cloglog"),
  weights = NULL,
  start = NULL,
  control = list()
)
\end{verbatim}
\end{Usage}
%
\begin{Arguments}
\begin{ldescription}
\item[\code{formula}] Formula with syntax: y \textasciitilde{} x + (terms | group)

\item[\code{data}] Data frame

\item[\code{link}] Link function: "logit" (default), "probit", or "cloglog"

\item[\code{weights}] Optional vector of case weights (one per observation)

\item[\code{start}] Optional starting values

\item[\code{control}] Control parameters for optimization
\end{ldescription}
\end{Arguments}
%
\begin{Details}
For binary response Y in \{0, 1\} or binomial response Y/n, the model is:

\strong{Logit link (default - logistic regression):}
\deqn{P(Y=1|x) = \frac{1}{1 + \exp(-x'\beta - Z'u)}}{}

\strong{Probit link:}
\deqn{P(Y=1|x) = \Phi(x'\beta + Z'u)}{}
where \eqn{\Phi}{} is the standard normal CDF.

\strong{Complementary log-log (cloglog) link:}
\deqn{P(Y=1|x) = 1 - \exp(-\exp(x'\beta + Z'u))}{}

The cloglog link is asymmetric and particularly useful for:
\begin{itemize}

\item{} Rare events (when baseline probability is low)
\item{} Survival analysis with discrete time
\item{} When modeling hazards that are proportional
\item{} Grouped survival data from an underlying Poisson process

\end{itemize}


Random effects u are assumed multivariate normal with mean 0.
\end{Details}
%
\begin{Value}
An object of class \code{gllamm\_binomial} with components:
\begin{ldescription}
\item[\code{coefficients}] List of fixed effects and random effect variances
\item[\code{logLik}] Log-likelihood at convergence
\item[\code{AIC}] Akaike Information Criterion
\item[\code{BIC}] Bayesian Information Criterion
\item[\code{convergence}] Convergence information
\item[\code{link}] Link function used
\item[\code{n\_obs}] Number of observations
\item[\code{fitted\_values}] Fitted probabilities
\item[\code{tmb\_obj}] TMB object for further inference
\item[\code{tmb\_opt}] Optimization result
\item[\code{tmb\_sdr}] Standard errors via sdreport
\end{ldescription}
\end{Value}
%
\begin{Note}
The recommended interface is \code{gllamm(formula, data, family = binomial(link))}.
This function is also available for direct use with the \code{link} argument.
\end{Note}
%
\begin{SeeAlso}
\code{\LinkA{gllamm}{gllamm}}, \code{\LinkA{binomial}{binomial}}, \code{\LinkA{ordinal}{ordinal}}
\end{SeeAlso}
%
\begin{Examples}
\begin{ExampleCode}
## Not run: 
# Simulate binary data
set.seed(123)
n_groups <- 20
n_per_group <- 10
data <- data.frame(
  group = rep(1:n_groups, each = n_per_group),
  x = rnorm(n_groups * n_per_group),
  y = rbinom(n_groups * n_per_group, 1, 0.5)
)

# Recommended: Use gllamm() with binomial() family
fit1 <- gllamm(y ~ x + (1 | group),
               data = data,
               family = binomial(link = "logit"))
summary(fit1)

# Probit link
fit2 <- gllamm(y ~ x + (1 | group),
               data = data,
               family = binomial(link = "probit"))

# Complementary log-log for rare events
data$rare_event <- rbinom(nrow(data), 1, 0.05)
fit3 <- gllamm(rare_event ~ x + (1 | group),
               data = data,
               family = binomial(link = "cloglog"))
summary(fit3)

# Alternative: Call fit_binomial() directly
fit4 <- fit_binomial(y ~ x + (1 | group),
                     data = data,
                     link = "logit")

## End(Not run)

\end{ExampleCode}
\end{Examples}
\HeaderA{fit\_eirt}{Fit Explanatory Item Response Theory Models}{fit.Rul.eirt}
%
\begin{Description}
Fit IRT models where item parameters are modeled as functions of item covariates.
For polytomous models, supports GRM (cumulative logit), PCM (adjacent-categories),
and GPCM (adjacent-categories with discrimination). PCM can include optional
threshold-level predictors via threshold\_formula.
\end{Description}
%
\begin{Usage}
\begin{verbatim}
fit_eirt(
  response_matrix,
  item_data,
  difficulty_formula = ~1,
  discrimination_formula = ~1,
  threshold_formula = NULL,
  person_data = NULL,
  random = NULL,
  weights = NULL,
  model = c("Rasch", "2PL", "GRM", "PCM", "GPCM"),
  item_residuals = TRUE,
  start = NULL,
  control = list()
)
\end{verbatim}
\end{Usage}
%
\begin{Arguments}
\begin{ldescription}
\item[\code{response\_matrix}] Matrix of item responses (persons x items)

\item[\code{item\_data}] Data frame of item-level covariates (must have n\_items rows)

\item[\code{difficulty\_formula}] Formula for item location/difficulty regression (e.g., \textasciitilde{} word\_freq).
Model: b\_i = W\_diff \%*\% gamma + epsilon\_b (if item\_residuals = TRUE)

\item[\code{discrimination\_formula}] Formula for discrimination regression (e.g., \textasciitilde{} item\_type).
Applies to 2PL, GRM, and GPCM models.
Model: log(a\_i) = W\_disc \%*\% delta + epsilon\_a (if item\_residuals = TRUE)

\item[\code{threshold\_formula}] Formula for threshold-specific regression (e.g., \textasciitilde{} abstractness).
Only used for PCM/GPCM models. When specified, enables threshold-difficulty regression
(Kim \& Wilson 2019 LPCM framework): delta\_im = b\_i + sum\_k xi\_km * x\_ik + e\_im

\item[\code{person\_data}] Optional data frame with person-level variables for multi-level models

\item[\code{random}] Optional random effects formula (e.g., \textasciitilde{} (1 | class))

\item[\code{weights}] Optional vector of case weights (fweights or pweights).
Length must equal number of persons. Default: all observations weighted equally.
Weights are expanded to observation-level (person-item) internally.

\item[\code{model}] IRT model type: "Rasch" or "2PL" (dichotomous), or
"GRM", "PCM", "GPCM" (polytomous)

\item[\code{item\_residuals}] Logical. If TRUE (default), includes item-specific residuals
(LLTM + error). If FALSE, uses pure LLTM where item parameters are exactly
predicted by covariates with no residuals.

\item[\code{start}] Optional starting values list

\item[\code{control}] Control parameters for nlminb optimization
\end{ldescription}
\end{Arguments}
%
\begin{Details}
**Dichotomous models:**
- **Rasch**: P(Y=1) = logit\textasciicircum{}(-1)(theta - b\_i), where b\_i = W\_diff \%*\% gamma [+ epsilon\_b]
- **2PL**: P(Y=1) = logit\textasciicircum{}(-1)(a\_i * (theta - b\_i)), where log(a\_i) = W\_disc \%*\% delta [+ epsilon\_a]

**Polytomous models (Kim \& Wilson 2019; De Boeck \& Wilson 2004 framework):**
- **GRM**: Cumulative logit with ordered thresholds. b\_i provides item location;
step\_param provides threshold spacing. Supports discrimination\_formula.
- **PCM**: Adjacent-categories logit, two-fold parameterization (MFRM approach):
delta\_im = b\_i + s\_im, where sum(s\_im) = 0 across steps for each item.
When threshold\_formula is specified, uses threshold regression:
delta\_im = b\_i + sum\_k xi\_km * x\_ik + e\_im (Kim \& Wilson LPCM framework)
- **GPCM**: As PCM with item-specific discrimination a\_i.
Supports both discrimination\_formula and threshold\_formula.

**Item residuals:** When item\_residuals = TRUE (default), item parameters include
residual terms epsilon\_b and epsilon\_a (LLTM + error). When FALSE, uses pure LLTM
where parameters are exactly determined by covariates.
\end{Details}
%
\begin{Value}
An object of class \code{gllamm\_eirt}
\end{Value}
%
\begin{Examples}
\begin{ExampleCode}
## Not run: 
# Dichotomous EIRT with item and discrimination predictors
item_chars <- data.frame(
  word_frequency = rnorm(20),
  item_length = rpois(20, 5),
  item_type = factor(sample(c("concrete", "abstract"), 20, replace = TRUE))
)
fit_2pl <- fit_eirt(responses, item_data = item_chars,
                    difficulty_formula = ~ word_frequency + item_length,
                    discrimination_formula = ~ item_type,
                    model = "2PL")

# Pure LLTM (no residuals)
fit_lltm <- fit_eirt(responses, item_data = item_chars,
                     difficulty_formula = ~ word_frequency,
                     model = "Rasch",
                     item_residuals = FALSE)

# Polytomous PCM (adjacent-categories logit, Rasch family)
fit_pcm <- fit_eirt(poly_responses, item_data = item_chars,
                    difficulty_formula = ~ abstractness,
                    model = "PCM")

# PCM with threshold predictors (LPCM framework)
fit_pcm_thresh <- fit_eirt(poly_responses, item_data = item_chars,
                           difficulty_formula = ~ abstractness,
                           threshold_formula = ~ cognitive_level,
                           model = "PCM")

# GPCM with all predictors
fit_gpcm <- fit_eirt(poly_responses, item_data = item_chars,
                     difficulty_formula = ~ word_frequency,
                     discrimination_formula = ~ item_type,
                     threshold_formula = ~ cognitive_level,
                     model = "GPCM")

## End(Not run)

\end{ExampleCode}
\end{Examples}
\HeaderA{fit\_irt}{Fit Item Response Theory Models}{fit.Rul.irt}
%
\begin{Description}
Fit dichotomous (Rasch, 2PL, 3PL) or polytomous (GRM, PCM, GPCM, NRM) IRT models.
Optionally include multi-level random effects for hierarchical or clustered data.
\end{Description}
%
\begin{Usage}
\begin{verbatim}
fit_irt(
  response_matrix,
  model = c("Rasch", "2PL", "3PL", "GRM", "PCM", "GPCM", "NRM"),
  person_data = NULL,
  random = NULL,
  weights = NULL,
  mc_items = NULL,
  method = c("auto", "em", "laplace"),
  quad_points = 61,
  se = FALSE,
  start = NULL,
  control = list()
)
\end{verbatim}
\end{Usage}
%
\begin{Arguments}
\begin{ldescription}
\item[\code{response\_matrix}] Matrix of item responses (persons x items).
For dichotomous models: coded 0/1.
For polytomous models: coded 1, 2, ..., K (ordered categories).

\item[\code{model}] Type of IRT model:
Dichotomous: "Rasch", "2PL", "3PL"
Polytomous: "GRM" (Graded Response), "PCM" (Partial Credit),
"GPCM" (Generalized Partial Credit), "NRM" (Nominal Response)

\item[\code{person\_data}] Optional data frame with person-level variables for multi-level models.
Must have one row per person (rows correspond to rows in response\_matrix).
Used to specify grouping variables in the random effects formula.

\item[\code{random}] Optional random effects formula using lme4-style syntax.
Examples: \code{\textasciitilde{} (1 | class)}, \code{\textasciitilde{} (1 | school/class)},
\code{\textasciitilde{} (1 | student) + (1 | time)}.
Requires person\_data to contain the grouping variables.

\item[\code{weights}] Optional vector of case weights. Length must equal number of persons.

\item[\code{mc\_items}] For 3PL model only: which items have guessing parameters.
Can be: NULL (default, all items have guessing), logical vector (length = n\_items),
or integer vector (indices of MC items). Non-MC items use 2PL likelihood (no guessing).

\item[\code{method}] Estimation method. "auto" (default) uses "em" for
single-level models and "laplace" whenever multi-level structure
(\code{random}) or standard errors (\code{se = TRUE}) require it.
"em" is Bock-Aitkin marginal maximum likelihood with fixed
Gauss-Hermite quadrature (the mirt/TAM algorithm; typically 20-50x
faster and evaluates the marginal likelihood more exactly than the
Laplace approximation). "laplace" is the TMB path, required for
multi-level models. EM person abilities are EAP scores; Laplace
abilities are posterior modes.

\item[\code{quad\_points}] Number of quadrature nodes for method = "em"
(default 61)

\item[\code{se}] Compute parameter standard errors via TMB::sdreport (default
FALSE, matching the default behavior of mirt; SE computation roughly
doubles the fitting time for large person samples)

\item[\code{start}] Optional starting values

\item[\code{control}] Control parameters for optimization
\end{ldescription}
\end{Arguments}
%
\begin{Value}
An object of class \code{gllamm\_irt}
\end{Value}
%
\begin{Examples}
\begin{ExampleCode}
## Not run: 
# Dichotomous example (Rasch)
set.seed(123)
n_persons <- 500
n_items <- 20
theta <- rnorm(n_persons, 0, 1)
difficulty <- rnorm(n_items, 0, 1)

# Generate binary responses
responses <- matrix(NA, n_persons, n_items)
for (i in 1:n_persons) {
  for (j in 1:n_items) {
    p <- plogis(theta[i] - difficulty[j])
    responses[i, j] <- rbinom(1, 1, p)
  }
}

# Fit Rasch model
fit_rasch <- fit_irt(responses, model = "Rasch")
summary(fit_rasch)

# Polytomous example (GRM)
# Generate 5-category responses
responses_poly <- matrix(NA, n_persons, n_items)
thresholds <- matrix(seq(-2, 2, length.out = 4), n_items, 4, byrow = TRUE)
for (i in 1:n_persons) {
  for (j in 1:n_items) {
    probs <- c(plogis(theta[i] - thresholds[j, 1]),
               diff(plogis(theta[i] - thresholds[j, ])),
               1 - plogis(theta[i] - thresholds[j, 4]))
    responses_poly[i, j] <- sample(1:5, 1, prob = probs)
  }
}

# Fit GRM model
fit_grm <- fit_irt(responses_poly, model = "GRM")
summary(fit_grm)

# 3PL with selective guessing (mixed MC and non-MC items)
# Assessment: 20 items, first 15 are MC, last 5 are open-ended
fit_3pl <- fit_irt(responses, model = "3PL", mc_items = 1:15)
# Only items 1-15 get guessing parameters
# Items 16-20 use 2PL likelihood (no guessing)

# Multi-level IRT: students nested in classes
person_data <- data.frame(
  person_id = 1:n_persons,
  class_id = rep(1:10, each = 50)
)
fit_multilevel <- fit_irt(responses, model = "2PL",
                           person_data = person_data,
                           random = ~ (1 | class_id))
# theta_i = theta_0i + u_class[class[i]]

## End(Not run)

\end{ExampleCode}
\end{Examples}
\HeaderA{fit\_irt\_dichotomous}{Internal function for dichotomous IRT models}{fit.Rul.irt.Rul.dichotomous}
\keyword{internal}{fit\_irt\_dichotomous}
%
\begin{Description}
Internal function for dichotomous IRT models
\end{Description}
%
\begin{Usage}
\begin{verbatim}
fit_irt_dichotomous(
  response_matrix,
  model,
  weights,
  mc_items,
  re_info,
  se,
  start,
  control
)
\end{verbatim}
\end{Usage}
\HeaderA{fit\_irt\_em}{MML-EM estimation for IRT models (Bock-Aitkin)}{fit.Rul.irt.Rul.em}
\keyword{internal}{fit\_irt\_em}
%
\begin{Description}
Marginal maximum likelihood via the EM algorithm with fixed Gauss-Hermite
quadrature over the ability distribution - the algorithm class used by
mirt and TAM. The E-step computes person-by-node posterior weights in one
matrix product; M-steps are independent small optimizations per item.
\end{Description}
%
\begin{Usage}
\begin{verbatim}
fit_irt_em(
  response_matrix,
  model,
  weights = NULL,
  mc_items = NULL,
  quad_points = 61,
  max_iter = 500,
  tol = 1e-04,
  control = list()
)
\end{verbatim}
\end{Usage}
%
\begin{Details}
Identification matches the Laplace path: 2PL/3PL/GRM/GPCM/NRM fix the
ability SD at 1; Rasch and PCM keep a free ability SD via the
common-slope equivalence (theta \textasciitilde{} N(0, sigma) with unit slope is
theta \textasciitilde{} N(0,1) with shared slope sigma; difficulties back-transform as
b = d * sigma).
\end{Details}
\HeaderA{fit\_irt\_polytomous}{Internal function for polytomous IRT models}{fit.Rul.irt.Rul.polytomous}
\keyword{internal}{fit\_irt\_polytomous}
%
\begin{Description}
Internal function for polytomous IRT models
\end{Description}
%
\begin{Usage}
\begin{verbatim}
fit_irt_polytomous(
  response_matrix,
  model,
  weights,
  re_info,
  se,
  start,
  control
)
\end{verbatim}
\end{Usage}
\HeaderA{fit\_lca}{Fit Latent Class Analysis Models}{fit.Rul.lca}
%
\begin{Description}
Fit latent class models to categorical data
\end{Description}
%
\begin{Usage}
\begin{verbatim}
fit_lca(
  formula,
  data = NULL,
  nclass = 2,
  weights = NULL,
  method = c("em", "tmb"),
  start = NULL,
  control = list()
)
\end{verbatim}
\end{Usage}
%
\begin{Arguments}
\begin{ldescription}
\item[\code{formula}] Formula specifying manifest variables (can be a matrix)

\item[\code{data}] Data frame containing variables

\item[\code{nclass}] Number of latent classes

\item[\code{weights}] Optional vector of case weights (one per observation)

\item[\code{method}] Estimation method: "em" (default; closed-form EM with
Ramsay acceleration, the poLCA algorithm) or "tmb" (direct
quasi-Newton on the marginal likelihood via TMB)

\item[\code{start}] Optional starting values

\item[\code{control}] Control parameters
\end{ldescription}
\end{Arguments}
%
\begin{Value}
An object of class \code{gllamm\_lca}
\end{Value}
%
\begin{Examples}
\begin{ExampleCode}
## Not run: 
# Simulate 2-class data
set.seed(123)
n <- 500

# Class 1: high probability of yes
class1_probs <- c(0.8, 0.7, 0.9, 0.75)
# Class 2: low probability of yes
class2_probs <- c(0.2, 0.3, 0.1, 0.25)

# Generate data
true_class <- sample(1:2, n, replace = TRUE, prob = c(0.6, 0.4))
data <- matrix(NA, n, 4)
for (i in 1:n) {
  probs <- if (true_class[i] == 1) class1_probs else class2_probs
  data[i, ] <- rbinom(4, 1, probs)
}
colnames(data) <- paste0("Item", 1:4)

# Fit 2-class model
fit <- fit_lca(data, nclass = 2)
summary(fit)

## End(Not run)

\end{ExampleCode}
\end{Examples}
\HeaderA{fit\_lca\_em}{EM estimation for latent class models}{fit.Rul.lca.Rul.em}
\keyword{internal}{fit\_lca\_em}
%
\begin{Description}
Classic EM for finite mixtures with conditionally independent indicators
(the poLCA algorithm). All M-steps are closed-form (weighted proportions
for binary/categorical indicators, weighted moments for gaussian ones),
so each iteration is one N x K posterior computation plus a handful of
cross-products - the E-step runs on BLAS matrix products.
\end{Description}
%
\begin{Usage}
\begin{verbatim}
fit_lca_em(
  Y,
  nclass,
  item_type,
  n_cats,
  weights = NULL,
  n_starts = 3,
  max_iter = 1000,
  tol = 1e-08
)
\end{verbatim}
\end{Usage}
\HeaderA{fit\_mixed}{Fit Joint Models for Mixed Response Types}{fit.Rul.mixed}
%
\begin{Description}
Fits a joint model for up to three outcomes of different types
(gaussian, binomial, poisson) measured on the same observations and
sharing a common random effect. The shared random effect induces
association between the outcomes.
\end{Description}
%
\begin{Usage}
\begin{verbatim}
fit_mixed(formulas, random, data, start = NULL, control = list())
\end{verbatim}
\end{Usage}
%
\begin{Arguments}
\begin{ldescription}
\item[\code{formulas}] Named list of up to three fixed-effects formulas, with
names among "gaussian", "binomial", "poisson", e.g.
\code{list(gaussian = y1 \textasciitilde{} x1, binomial = y2 \textasciitilde{} x2)}.

\item[\code{random}] One-sided random-intercept formula, e.g. \code{\textasciitilde{} (1 | group)}

\item[\code{data}] Data frame containing all outcome and covariate variables;
rows must be complete for every supplied outcome

\item[\code{start}] Optional starting values

\item[\code{control}] Optimization control list
\end{ldescription}
\end{Arguments}
%
\begin{Details}
The shared random effect enters every outcome's linear predictor with
loading 1 (a common-intercept joint model). Outcome-specific loadings
are not yet supported.
\end{Details}
%
\begin{Value}
An object of class \code{gllamm\_mixed}
\end{Value}
%
\begin{Examples}
\begin{ExampleCode}
## Not run: 
fit <- fit_mixed(
  formulas = list(gaussian = biomarker ~ age,
                  binomial = event ~ age + treatment),
  random = ~ (1 | patient),
  data = d)

## End(Not run)

\end{ExampleCode}
\end{Examples}
\HeaderA{fit\_multinomial}{Fit Multinomial Regression Models with Random Effects}{fit.Rul.multinomial}
%
\begin{Description}
Fit baseline category logit models for nominal (unordered) responses
\end{Description}
%
\begin{Usage}
\begin{verbatim}
fit_multinomial(
  formula,
  data,
  reference = NULL,
  start = NULL,
  control = list()
)
\end{verbatim}
\end{Usage}
%
\begin{Arguments}
\begin{ldescription}
\item[\code{formula}] Formula with syntax: y \textasciitilde{} x + (terms | group)

\item[\code{data}] Data frame

\item[\code{reference}] Reference category (default: first level)

\item[\code{start}] Optional starting values

\item[\code{control}] Control parameters
\end{ldescription}
\end{Arguments}
%
\begin{Details}
For nominal response Y with K categories, using baseline category logit:

\deqn{P(Y = k | x) = \frac{\exp(x'\beta_k)}{1 + \sum_{j=1}^{K-1} \exp(x'\beta_j)}}{}

where category 0 is the reference with \eqn{\beta_0 = 0}{}.
\end{Details}
%
\begin{Value}
An object of class \code{gllamm\_multinomial}
\end{Value}
%
\begin{Examples}
\begin{ExampleCode}
## Not run: 
# Simulate multinomial data
data$choice <- factor(sample(c("A", "B", "C"), 100, replace = TRUE))

# Fit multinomial model
fit <- fit_multinomial(choice ~ price + quality + (1 | person),
                       data = data)
summary(fit)

## End(Not run)

\end{ExampleCode}
\end{Examples}
\HeaderA{fit\_npml}{Fit Two-Level GLMMs by Nonparametric Maximum Likelihood (NPML)}{fit.Rul.npml}
%
\begin{Description}
Replaces the normal random-intercept distribution with a discrete
distribution on \code{k} estimated mass points (locations and masses).
The marginal likelihood is an exact finite mixture - no integral
approximation - making NPML robust to misspecification of the
random-effects distribution.
\end{Description}
%
\begin{Usage}
\begin{verbatim}
fit_npml(
  formula,
  data,
  k = 2,
  family = stats::gaussian(),
  weights = NULL,
  n_starts = 3,
  start = NULL,
  control = list()
)
\end{verbatim}
\end{Usage}
%
\begin{Arguments}
\begin{ldescription}
\item[\code{formula}] Model formula \code{y \textasciitilde{} x + (1 | group)}

\item[\code{data}] Data frame

\item[\code{k}] Number of mass points (default 2)

\item[\code{family}] gaussian(), binomial(), or poisson()

\item[\code{weights}] Optional observation weights

\item[\code{n\_starts}] Random restarts (mixtures are prone to local optima)

\item[\code{start}] Optional starting values

\item[\code{control}] Optimization control list
\end{ldescription}
\end{Arguments}
%
\begin{Details}
The fixed-effects design drops its intercept: the mass-point locations
play the role of component intercepts (the usual NPML identification,
as in \pkg{npmlreg}). Masses are softmax-parameterized.
\end{Details}
%
\begin{Value}
An object of class \code{gllamm\_npml}
\end{Value}
%
\begin{Examples}
\begin{ExampleCode}
## Not run: 
fit <- fit_npml(y ~ x + (1 | group), data = d, k = 3,
                family = stats::binomial())

## End(Not run)

\end{ExampleCode}
\end{Examples}
\HeaderA{fit\_ordinal}{Fit Ordinal Regression Models with Random Effects}{fit.Rul.ordinal}
%
\begin{Description}
Fit proportional odds or cumulative probit models for ordinal responses.
This function can be called directly or through \code{gllamm()} with
\code{family = ordinal()}. The \code{gllamm()} interface is recommended
for consistency with other model types.
\end{Description}
%
\begin{Usage}
\begin{verbatim}
fit_ordinal(
  formula,
  data,
  link = c("logit", "probit", "acl", "crl_forward", "crl_backward", "ppo"),
  weights = NULL,
  start = NULL,
  control = list()
)
\end{verbatim}
\end{Usage}
%
\begin{Arguments}
\begin{ldescription}
\item[\code{formula}] Formula with syntax: y \textasciitilde{} x + (terms | group)

\item[\code{data}] Data frame

\item[\code{link}] Link function: "logit" (proportional odds) or "probit"

\item[\code{weights}] Optional vector of case weights (one per observation)

\item[\code{start}] Optional starting values

\item[\code{control}] Control parameters
\end{ldescription}
\end{Arguments}
%
\begin{Details}
For ordinal response Y with K categories (1, 2, ..., K), the model is:

Proportional odds (logit link):
\deqn{P(Y \le k | x) = \frac{1}{1 + \exp(-(\tau_k - x'\beta))}}{}

Cumulative probit:
\deqn{P(Y \le k | x) = \Phi(\tau_k - x'\beta)}{}

where \eqn{\tau_1 < \tau_2 < ... < \tau_{K-1}}{} are threshold parameters.
\end{Details}
%
\begin{Value}
An object of class \code{gllamm\_ordinal}
\end{Value}
%
\begin{Note}
The recommended interface is \code{gllamm(formula, data, family = ordinal(link))}.
This function is also available for direct use with the \code{link} argument.
\end{Note}
%
\begin{Examples}
\begin{ExampleCode}
## Not run: 
# Simulate ordinal data
data$satisfaction <- factor(sample(1:5, 100, replace = TRUE),
                            ordered = TRUE,
                            levels = 1:5)

# Recommended: Use gllamm() with ordinal() family
fit1 <- gllamm(satisfaction ~ age + (1 | clinic),
               data = data,
               family = ordinal(link = "logit"))
summary(fit1)

# Alternative: Call fit_ordinal() directly
fit2 <- fit_ordinal(satisfaction ~ age + (1 | clinic),
                    data = data,
                    link = "logit")
summary(fit2)

## End(Not run)

\end{ExampleCode}
\end{Examples}
\HeaderA{fit\_rank}{Fit Rank-Ordered Logit Models with Taste Heterogeneity}{fit.Rul.rank}
%
\begin{Description}
Fits an exploded (rank-ordered) logit model for ranking data, with an
optional group-level random coefficient on one alternative attribute
(a "taste shifter"). A constant-within-case random intercept cancels in
ranking likelihoods, so heterogeneity must attach to an attribute that
varies across alternatives.
\end{Description}
%
\begin{Usage}
\begin{verbatim}
fit_rank(
  formula,
  case,
  data,
  random = NULL,
  weights = NULL,
  start = NULL,
  control = list()
)
\end{verbatim}
\end{Usage}
%
\begin{Arguments}
\begin{ldescription}
\item[\code{formula}] Model formula \code{rank \textasciitilde{} x1 + x2} where the response is
the rank of each alternative within its case (1 = most preferred).
Unranked alternatives may be coded \code{NA}; they remain in every
choice set without contributing a stage.

\item[\code{case}] One-sided formula naming the case identifier, e.g. \code{\textasciitilde{} id}

\item[\code{data}] Data frame in long format (one row per alternative per case)

\item[\code{random}] Optional one-sided formula \code{\textasciitilde{} (0 + attribute | group)}
giving the attribute carrying the random coefficient and the grouping
variable. NULL fits a fixed-effects rank-ordered logit (the random SD
is fixed near zero).

\item[\code{weights}] Optional one weight per case (matched via the case id)

\item[\code{start}] Optional starting values

\item[\code{control}] Optimization control list
\end{ldescription}
\end{Arguments}
%
\begin{Value}
An object of class \code{gllamm\_rank}
\end{Value}
%
\begin{Examples}
\begin{ExampleCode}
## Not run: 
fit <- fit_rank(rank ~ price + quality, case = ~ subject,
                random = ~ (0 + price | region), data = d)

## End(Not run)

\end{ExampleCode}
\end{Examples}
\HeaderA{fit\_sem}{Fit Structural Equation Models with Latent Variables}{fit.Rul.sem}
%
\begin{Description}
Fits a SEM with continuous indicators: a measurement model defining each
latent variable by its indicators (marker-variable identification: the
first indicator's loading is fixed at 1) and an optional recursive
structural model of regressions among latent variables.
\end{Description}
%
\begin{Usage}
\begin{verbatim}
fit_sem(
  measurement,
  structural = NULL,
  data,
  method = c("ml", "laplace"),
  start = NULL,
  control = list()
)
\end{verbatim}
\end{Usage}
%
\begin{Arguments}
\begin{ldescription}
\item[\code{measurement}] Named list of one-sided formulas defining each latent
variable, e.g. \code{list(ability = \textasciitilde{} y1 + y2 + y3,
motivation = \textasciitilde{} y4 + y5 + y6)}.

\item[\code{structural}] Optional list of formulas among latent variables, e.g.
\code{list(motivation \textasciitilde{} ability)}. Must be recursive (no cycles).

\item[\code{data}] Data frame with the indicator variables

\item[\code{method}] Estimation method: "ml" (default; Wishart maximum
likelihood on the sample covariance matrix, the lavaan/LISREL
approach - fitting cost independent of N) or "laplace" (full-data
TMB Laplace estimation). Both give the same ML estimates for
complete data.

\item[\code{start}] Optional starting values (laplace method only)

\item[\code{control}] Optimization control list
\end{ldescription}
\end{Arguments}
%
\begin{Value}
An object of class \code{gllamm\_sem}
\end{Value}
%
\begin{Examples}
\begin{ExampleCode}
## Not run: 
fit <- fit_sem(
  measurement = list(f1 = ~ x1 + x2 + x3, f2 = ~ y1 + y2 + y3),
  structural = list(f2 ~ f1),
  data = d)

## End(Not run)

\end{ExampleCode}
\end{Examples}
\HeaderA{fit\_sem\_ml}{Covariance-based ML estimation for SEM (Wishart likelihood)}{fit.Rul.sem.Rul.ml}
\keyword{internal}{fit\_sem\_ml}
%
\begin{Description}
Fits the measurement + recursive structural model from the sample
covariance matrix by minimizing the ML discrepancy
F = log|Sigma(theta)| + tr(S Sigma\textasciicircum{}-1) - log|S| - p,
the lavaan/LISREL approach. The data reduce to S and the mean vector, so
fitting cost is independent of N. For complete data this yields the same
ML estimates as the full-data marginal likelihood.
\end{Description}
%
\begin{Usage}
\begin{verbatim}
fit_sem_ml(Y, lambda_pattern, beta_pattern, control = list())
\end{verbatim}
\end{Usage}
\HeaderA{fit\_statistics}{Model Fit Statistics}{fit.Rul.statistics}
%
\begin{Description}
Extract comprehensive fit statistics from GLLAMM models
\end{Description}
\HeaderA{fit\_survival}{Fit Parametric Survival Models with Random Effects (Frailty)}{fit.Rul.survival}
%
\begin{Description}
Fits exponential or Weibull proportional-hazards models with log-normal
frailty (normally distributed random effects on the log-hazard scale),
supporting right censoring.
\end{Description}
%
\begin{Usage}
\begin{verbatim}
fit_survival(
  formula,
  data,
  distribution = c("weibull", "exponential"),
  weights = NULL,
  start = NULL,
  control = list()
)
\end{verbatim}
\end{Usage}
%
\begin{Arguments}
\begin{ldescription}
\item[\code{formula}] Formula of the form \code{Surv(time, event) \textasciitilde{} x + (1 | group)}.
The left-hand side names the time and event (1 = event, 0 = censored)
variables; the \pkg{survival} package is not required.

\item[\code{data}] Data frame

\item[\code{distribution}] "weibull" (default) or "exponential"

\item[\code{weights}] Optional case weights

\item[\code{start}] Optional starting values

\item[\code{control}] Optimization control list
\end{ldescription}
\end{Arguments}
%
\begin{Details}
The cumulative hazard is \eqn{H(t \mid x, u) = (\lambda t)^{shape}}{} with
\eqn{\lambda = \exp(x'\beta + z'u)}{}; the exponential model fixes
shape = 1. Under this parameterization the accelerated-failure-time
coefficients of \code{survival::survreg} correspond to \eqn{-\beta}{} and
its scale to \eqn{1/shape}{}.

The exponential frailty model is likelihood-equivalent to a Poisson GLMM
on the event indicator with offset \eqn{\log(t)}{}, which is used in the
package validation suite.
\end{Details}
%
\begin{Value}
An object of class \code{gllamm\_survival}
\end{Value}
%
\begin{Examples}
\begin{ExampleCode}
## Not run: 
fit <- fit_survival(Surv(time, status) ~ age + (1 | center),
                    data = d, distribution = "weibull")

## End(Not run)

\end{ExampleCode}
\end{Examples}
\HeaderA{fit\_tmb\_gllamm}{Interface to TMB for GLLAMM models}{fit.Rul.tmb.Rul.gllamm}
\keyword{internal}{fit\_tmb\_gllamm}
%
\begin{Description}
Prepares data and parameters for TMB, compiles and runs optimization
\end{Description}
%
\begin{Usage}
\begin{verbatim}
fit_tmb_gllamm(
  model_data,
  family,
  start_params = NULL,
  control = list(),
  weights = NULL
)
\end{verbatim}
\end{Usage}
%
\begin{Arguments}
\begin{ldescription}
\item[\code{model\_data}] List from make\_model\_matrices()

\item[\code{family}] GLM family object

\item[\code{start\_params}] Optional starting values

\item[\code{control}] Control parameters for optimization
\end{ldescription}
\end{Arguments}
%
\begin{Value}
TMB fit object with parameter estimates
\end{Value}
\HeaderA{fit\_tmb\_gllamm\_aghq}{Fit a two-level GLMM by adaptive Gauss-Hermite quadrature}{fit.Rul.tmb.Rul.gllamm.Rul.aghq}
\keyword{internal}{fit\_tmb\_gllamm\_aghq}
%
\begin{Description}
R driver for the glmm\_aghq TMB objective: alternates parameter
optimization with updates of the per-group adaptation centers and scales
(posterior modes and curvatures), the classic adaptive quadrature scheme.
\end{Description}
%
\begin{Usage}
\begin{verbatim}
fit_tmb_gllamm_aghq(
  model_data,
  family,
  random_terms,
  k = 15,
  start_params = NULL,
  control = list(),
  weights = NULL,
  max_adapt = 5,
  adapt_tol = 1e-04
)
\end{verbatim}
\end{Usage}
\HeaderA{fit\_tmb\_gllamm\_multi}{GLMM engine for multiple random-effects terms (crossed/nested)}{fit.Rul.tmb.Rul.gllamm.Rul.multi}
\keyword{internal}{fit\_tmb\_gllamm\_multi}
%
\begin{Description}
Builds the lme4-style combined sparse Z mapping the full term-major
random-effects vector to observations, and fits via the glmm\_multi
TMB model.
\end{Description}
%
\begin{Usage}
\begin{verbatim}
fit_tmb_gllamm_multi(
  model_data,
  family,
  random_terms,
  start_params = NULL,
  control = list(),
  weights = NULL
)
\end{verbatim}
\end{Usage}
\HeaderA{fit\_tmb\_gllamm\_v2}{Enhanced interface to TMB for GLLAMMR models}{fit.Rul.tmb.Rul.gllamm.Rul.v2}
\keyword{internal}{fit\_tmb\_gllamm\_v2}
%
\begin{Description}
Supports random intercepts, random slopes, and multiple GLM families
\end{Description}
%
\begin{Usage}
\begin{verbatim}
fit_tmb_gllamm_v2(
  model_data,
  family,
  random_terms,
  start_params = NULL,
  control = list(),
  weights = NULL
)
\end{verbatim}
\end{Usage}
%
\begin{Arguments}
\begin{ldescription}
\item[\code{model\_data}] List from make\_model\_matrices()

\item[\code{family}] GLM family object

\item[\code{random\_terms}] Parsed random effects terms

\item[\code{start\_params}] Optional starting values

\item[\code{control}] Control parameters for optimization
\end{ldescription}
\end{Arguments}
%
\begin{Value}
TMB fit object with parameter estimates
\end{Value}
\HeaderA{fit\_tmb\_objective\_only}{Build a TMB objective (no optimization) for one cluster}{fit.Rul.tmb.Rul.objective.Rul.only}
\keyword{internal}{fit\_tmb\_objective\_only}
%
\begin{Description}
Mirrors the single-term v2 engine's data/parameter construction so the
cluster objective is evaluated under exactly the same model.
\end{Description}
%
\begin{Usage}
\begin{verbatim}
fit_tmb_objective_only(model_data, family, random_terms)
\end{verbatim}
\end{Usage}
\HeaderA{fixef}{Generic fixef}{fixef}
%
\begin{Description}
Generic fixef
\end{Description}
%
\begin{Usage}
\begin{verbatim}
fixef(object, ...)
\end{verbatim}
\end{Usage}
%
\begin{Arguments}
\begin{ldescription}
\item[\code{object}] A fitted model object

\item[\code{...}] Additional arguments
\end{ldescription}
\end{Arguments}
\HeaderA{fixef.gllamm}{Extract fixed effects}{fixef.gllamm}
%
\begin{Description}
Extract fixed effects
\end{Description}
%
\begin{Usage}
\begin{verbatim}
## S3 method for class 'gllamm'
fixef(object, ...)
\end{verbatim}
\end{Usage}
%
\begin{Arguments}
\begin{ldescription}
\item[\code{object}] A gllamm object

\item[\code{...}] Additional arguments (ignored)
\end{ldescription}
\end{Arguments}
%
\begin{Value}
Named vector of fixed effects coefficients
\end{Value}
\HeaderA{gauss\_hermite}{Gauss-Hermite nodes and weights (Golub-Welsch)}{gauss.Rul.hermite}
\keyword{internal}{gauss\_hermite}
%
\begin{Description}
Gauss-Hermite nodes and weights (Golub-Welsch)
\end{Description}
%
\begin{Usage}
\begin{verbatim}
gauss_hermite(n)
\end{verbatim}
\end{Usage}
\HeaderA{get\_inverse\_link}{Get inverse link function for a family}{get.Rul.inverse.Rul.link}
\keyword{internal}{get\_inverse\_link}
%
\begin{Description}
Get inverse link function for a family
\end{Description}
%
\begin{Usage}
\begin{verbatim}
get_inverse_link(family)
\end{verbatim}
\end{Usage}
%
\begin{Arguments}
\begin{ldescription}
\item[\code{family}] Family object (standard R family or custom)
\end{ldescription}
\end{Arguments}
%
\begin{Value}
Inverse link function
\end{Value}
\HeaderA{gllamm}{Fit Generalized Linear Latent and Mixed Models}{gllamm}
%
\begin{Description}
Main function for fitting GLLAMM models. Supports multilevel generalized
linear models, factor models, IRT models, latent class models, and more.
\end{Description}
%
\begin{Usage}
\begin{verbatim}
gllamm(
  formula,
  data,
  family = gaussian(),
  weights = NULL,
  random = NULL,
  integration = NULL,
  start = NULL,
  control = list(),
  ...
)
\end{verbatim}
\end{Usage}
%
\begin{Arguments}
\begin{ldescription}
\item[\code{formula}] A two-sided formula with syntax: \code{y \textasciitilde{} x + (terms | group)}.
The right-hand side contains fixed effects and random effects specifications.
Random effects are specified using \code{(term | group)} for correlated
random effects or \code{(term || group)} for uncorrelated random effects.
Nested random effects can be specified as \code{(1 | level1/level2)}.

\item[\code{data}] A data frame containing the variables in the formula.

\item[\code{family}] A GLM family object specifying the error distribution and link
function. Supports:
\begin{itemize}

\item{} \code{gaussian()} - Normal distribution (default)
\item{} \code{binomial()} - Binary/binomial outcomes
\item{} \code{poisson()} - Count data
\item{} \code{ordinal()} - Ordered categorical responses (see \code{?ordinal})

\end{itemize}


\item[\code{weights}] Optional case weights: a numeric vector of observation
(level-1) weights, or a list with elements \code{level1} and/or
\code{level2} for survey designs with weights at both levels.

\item[\code{random}] For matrix-response families (\code{irt()}, \code{lca()}):
an optional person-level random-effects formula such as
\code{\textasciitilde{} (1 | class)}.

\item[\code{integration}] Optional integration specification; \code{aghq(k)}
selects adaptive Gauss-Hermite quadrature with \code{k} nodes for
two-level random-intercept models (default is the Laplace
approximation).

\item[\code{start}] Optional named list of starting values for parameters.

\item[\code{control}] A list of control parameters for the optimization algorithm:
\begin{description}

\item[eval.max] Maximum number of function evaluations (default: 1000)
\item[iter.max] Maximum number of iterations (default: 500)
\item[trace] Integer controlling printed output (default: 0)

\end{description}


\item[\code{...}] Additional arguments (reserved for future use).
\end{ldescription}
\end{Arguments}
%
\begin{Details}
GLLAMMR uses Template Model Builder (TMB) for efficient computation via
automatic differentiation. The random effects are integrated out using
Laplace approximation, providing fast and accurate inference.

The formula syntax follows lme4 conventions:
\begin{itemize}

\item{} \code{(1 | group)} - Random intercept for group
\item{} \code{(x | group)} - Random intercept and slope for x
\item{} \code{(x || group)} - Uncorrelated random intercept and slope
\item{} \code{(1 | level1/level2)} - Nested random effects
\item{} \code{(1 | group1) + (1 | group2)} - Crossed random effects

\end{itemize}

\end{Details}
%
\begin{Value}
An object of class \code{gllamm} with components:
\begin{ldescription}
\item[\code{coefficients}] List with \code{fixed} (fixed effects coefficients)
and \code{random\_var} (random effects variance components)
\item[\code{vcov}] List with \code{fixed} (variance-covariance matrix of fixed
effects) and \code{all} (full variance-covariance matrix)
\item[\code{random\_effects}] List of random effects predictions by group
\item[\code{fitted.values}] Fitted values on the response scale
\item[\code{residuals}] Response residuals
\item[\code{y}] Response vector
\item[\code{X}] Fixed effects design matrix
\item[\code{logLik}] Log-likelihood at convergence
\item[\code{AIC}] Akaike Information Criterion
\item[\code{BIC}] Bayesian Information Criterion
\item[\code{n\_obs}] Number of observations
\item[\code{n\_params}] Number of parameters
\item[\code{n\_groups}] Number of groups for each random effects term
\item[\code{convergence}] List with convergence information
\item[\code{call}] The matched call
\item[\code{formula}] The model formula
\item[\code{family}] The GLM family
\item[\code{data}] The data frame (if requested)
\item[\code{random\_terms}] List of random effects specifications
\end{ldescription}
\end{Value}
%
\begin{References}
Rabe-Hesketh, S., Skrondal, A., \& Pickles, A. (2004). GLLAMM Manual.
U.C. Berkeley Division of Biostatistics Working Paper Series.

Skrondal, A., \& Rabe-Hesketh, S. (2004). Generalized Latent Variable
Modeling: Multilevel, Longitudinal, and Structural Equation Models.
Chapman \& Hall/CRC.
\end{References}
%
\begin{Examples}
\begin{ExampleCode}
## Not run: 
# Basic random intercept model
data(sleepstudy, package = "lme4")
fit1 <- gllamm(Reaction ~ Days + (1 | Subject),
               data = sleepstudy)
summary(fit1)

# Random intercept and slope
fit2 <- gllamm(Reaction ~ Days + (Days | Subject),
               data = sleepstudy)
summary(fit2)

# Extract components
fixef(fit2)        # Fixed effects
ranef(fit2)        # Random effects
VarCorr(fit2)      # Variance components
fitted(fit2)       # Fitted values
residuals(fit2)    # Residuals

# Ordinal regression (proportional odds)
data$satisfaction <- ordered(sample(1:5, nrow(data), replace = TRUE))
fit3 <- gllamm(satisfaction ~ x + (1 | group),
               data = data,
               family = ordinal(link = "logit"))

# Ordinal with adjacent category logit
fit4 <- gllamm(satisfaction ~ x + (1 | group),
               data = data,
               family = ordinal(link = "acl"))

## End(Not run)

\end{ExampleCode}
\end{Examples}
\HeaderA{gllamm-class}{GLLAMM model object class}{gllamm.Rdash.class}
\aliasA{coef.gllamm}{gllamm-class}{coef.gllamm}
\aliasA{fitted.gllamm}{gllamm-class}{fitted.gllamm}
\aliasA{logLik.gllamm}{gllamm-class}{logLik.gllamm}
\aliasA{print.gllamm}{gllamm-class}{print.gllamm}
\aliasA{residuals.gllamm}{gllamm-class}{residuals.gllamm}
\aliasA{summary.gllamm}{gllamm-class}{summary.gllamm}
\aliasA{vcov.gllamm}{gllamm-class}{vcov.gllamm}
%
\begin{Description}
S3 class for fitted GLLAMM models
\end{Description}
%
\begin{Usage}
\begin{verbatim}
## S3 method for class 'gllamm'
print(x, ...)

## S3 method for class 'gllamm'
summary(object, ...)

## S3 method for class 'gllamm'
coef(object, ...)

## S3 method for class 'gllamm'
vcov(object, which = "fixed", type = c("model", "sandwich"), ...)

## S3 method for class 'gllamm'
logLik(object, ...)

## S3 method for class 'gllamm'
fitted(object, ...)

## S3 method for class 'gllamm'
residuals(object, type = c("response", "pearson", "deviance"), ...)
\end{verbatim}
\end{Usage}
%
\begin{Arguments}
\begin{ldescription}
\item[\code{x}] A gllamm object

\item[\code{...}] Additional arguments

\item[\code{object}] A gllamm object

\item[\code{which}] Which covariance block to return: "fixed" (default) or "all"

\item[\code{type}] Covariance type: "model" (default) or "sandwich"
(cluster-robust)
\end{ldescription}
\end{Arguments}
\HeaderA{gllammr\_validate}{Cross-package validation of GLLAMMR estimates}{gllammr.Rul.validate}
%
\begin{Description}
Fits canonical benchmark datasets with GLLAMMR and with established
reference packages, and reports the agreement. Reference packages that use
the same Laplace approximation (lme4 with nAGQ = 1, ordinal::clmm) should
agree to numerical precision; packages using different integration schemes
(mirt, ltm EM quadrature) agree within small tolerances.
\end{Description}
%
\begin{Usage}
\begin{verbatim}
gllammr_validate(
  cases = "all",
  scale = c("standard", "large", "all"),
  verbose = TRUE
)
\end{verbatim}
\end{Usage}
%
\begin{Arguments}
\begin{ldescription}
\item[\code{cases}] Character vector of case names to run, or "all" (default).
Available: "gaussian\_sleepstudy", "binomial\_toenail",
"poisson\_grouseticks", "ordinal\_wine", "rasch\_lsat", "twopl\_simulated",
"lca\_carcinoma", "grm\_science", "gamma\_simulated",
"survival\_exponential", "sem\_lavaan", "lca\_polytomous",
"npml\_binomial", "aghq\_binomial", "twopl\_lsat\_em".

\item[\code{scale}] "standard" (default) runs the canonical-dataset cases;
"large" runs the large-scale tier (n in the tens of thousands, long
item batteries - sizes where quadrature grids and tolerances can fail
silently); "all" runs both.

\item[\code{verbose}] Print progress messages (default TRUE)
\end{ldescription}
\end{Arguments}
%
\begin{Details}
All reference packages are Suggests; cases whose reference package is not
installed are skipped.
\end{Details}
%
\begin{Value}
Data frame with one row per compared statistic: case, statistic,
GLLAMMR value, reference value, absolute and relative difference,
tolerance, and pass/fail.
\end{Value}
%
\begin{Examples}
\begin{ExampleCode}

if (requireNamespace("lme4", quietly = TRUE)) {
  gllammr_validate(cases = "gaussian_sleepstudy")
}


\end{ExampleCode}
\end{Examples}
\HeaderA{gof.gllamm}{Goodness-of-fit tests for GLLAMM}{gof.gllamm}
%
\begin{Description}
Goodness-of-fit tests for GLLAMM
\end{Description}
%
\begin{Usage}
\begin{verbatim}
gof.gllamm(object, ...)
\end{verbatim}
\end{Usage}
%
\begin{Arguments}
\begin{ldescription}
\item[\code{object}] A gllamm object

\item[\code{...}] Additional arguments
\end{ldescription}
\end{Arguments}
\HeaderA{icc}{Compute Intraclass Correlation Coefficients}{icc}
%
\begin{Description}
Calculate ICCs for multi-level IRT models
\end{Description}
%
\begin{Usage}
\begin{verbatim}
icc(x, ...)
\end{verbatim}
\end{Usage}
%
\begin{Arguments}
\begin{ldescription}
\item[\code{x}] A fitted multi-level IRT model

\item[\code{...}] Additional arguments (not used; methods may add \code{level})
\end{ldescription}
\end{Arguments}
%
\begin{Value}
A named vector of ICCs, or a single ICC if level specified
\end{Value}
%
\begin{Examples}
\begin{ExampleCode}
## Not run: 
# All ICCs
icc(fit)

# Specific level
icc(fit, level = "class")

## End(Not run)

\end{ExampleCode}
\end{Examples}
\HeaderA{icc.gllamm}{Variance decomposition (ICC)}{icc.gllamm}
%
\begin{Description}
Variance decomposition (ICC)
\end{Description}
%
\begin{Usage}
\begin{verbatim}
## S3 method for class 'gllamm'
icc(x, quiet = FALSE, ...)
\end{verbatim}
\end{Usage}
%
\begin{Arguments}
\begin{ldescription}
\item[\code{x}] A gllamm object

\item[\code{quiet}] Suppress printed output and messages (default: FALSE)

\item[\code{...}] Additional arguments
\end{ldescription}
\end{Arguments}
\HeaderA{influence.gllamm}{Influence diagnostics for GLLAMM}{influence.gllamm}
%
\begin{Description}
Influence diagnostics for GLLAMM
\end{Description}
%
\begin{Usage}
\begin{verbatim}
## S3 method for class 'gllamm'
influence(model, ...)
\end{verbatim}
\end{Usage}
%
\begin{Arguments}
\begin{ldescription}
\item[\code{model}] A gllamm object

\item[\code{...}] Additional arguments
\end{ldescription}
\end{Arguments}
%
\begin{Value}
Data frame with influence measures
\end{Value}
\HeaderA{irt}{IRT Family for Item Response Theory Models}{irt}
%
\begin{Description}
Create a family object for fitting IRT models through the unified
\code{gllamm()} interface. The response is a persons x items matrix
passed as the first argument of \code{gllamm()}.
\end{Description}
%
\begin{Usage}
\begin{verbatim}
irt(
  model = c("Rasch", "2PL", "3PL", "GRM", "PCM", "GPCM", "NRM"),
  mc_items = NULL
)
\end{verbatim}
\end{Usage}
%
\begin{Arguments}
\begin{ldescription}
\item[\code{model}] IRT model type: "Rasch", "2PL", "3PL" (dichotomous) or
"GRM", "PCM", "GPCM", "NRM" (polytomous)

\item[\code{mc\_items}] For 3PL only: which items have guessing parameters
(NULL = all; logical or integer index vector)
\end{ldescription}
\end{Arguments}
%
\begin{Value}
A family object of class \code{irt\_family}
\end{Value}
%
\begin{Examples}
\begin{ExampleCode}
## Not run: 
fit <- gllamm(response_matrix, family = irt("2PL"))
# Multi-level IRT: persons nested in classes
fit_ml <- gllamm(response_matrix, data = person_data,
                 family = irt("Rasch"), random = ~ (1 | class))

## End(Not run)

\end{ExampleCode}
\end{Examples}
\HeaderA{irt\_category\_probs}{Category response probabilities for polytomous IRT models}{irt.Rul.category.Rul.probs}
\keyword{internal}{irt\_category\_probs}
%
\begin{Description}
Single source of truth for the category probability math, mirroring the
likelihood in the TMB template (gllamm\_irt\_poly.hpp / gllamm\_eirt.hpp).
Used by plotting, DIF displays, and marginal predictions so the formulas
cannot drift apart.
\end{Description}
%
\begin{Usage}
\begin{verbatim}
irt_category_probs(model, theta, thresholds, discrimination = 1)
\end{verbatim}
\end{Usage}
%
\begin{Arguments}
\begin{ldescription}
\item[\code{model}] One of "GRM", "PCM", "GPCM", "NRM"

\item[\code{theta}] Numeric vector of ability values

\item[\code{thresholds}] Numeric vector of item parameters: ordered thresholds
(GRM), free step difficulties (PCM/GPCM), or category intercepts (NRM,
reference category omitted)

\item[\code{discrimination}] Item discrimination (ignored for PCM)
\end{ldescription}
\end{Arguments}
%
\begin{Value}
Matrix [length(theta) x K] of category probabilities; rows sum to 1
\end{Value}
\HeaderA{lca}{Latent Class Family for Finite Mixture Models}{lca}
%
\begin{Description}
Create a family object for fitting latent class models through the unified
\code{gllamm()} interface. The response is a matrix of binary manifest
variables passed as the first argument of \code{gllamm()}.
\end{Description}
%
\begin{Usage}
\begin{verbatim}
lca(nclass = 2)
\end{verbatim}
\end{Usage}
%
\begin{Arguments}
\begin{ldescription}
\item[\code{nclass}] Number of latent classes (default 2)
\end{ldescription}
\end{Arguments}
%
\begin{Value}
A family object of class \code{lca\_family}
\end{Value}
%
\begin{Examples}
\begin{ExampleCode}
## Not run: 
fit <- gllamm(indicator_matrix, family = lca(nclass = 3))

## End(Not run)

\end{ExampleCode}
\end{Examples}
\HeaderA{make\_model\_matrices}{Extract model matrices}{make.Rul.model.Rul.matrices}
\keyword{internal}{make\_model\_matrices}
%
\begin{Description}
Create design matrices for fixed and random effects
\end{Description}
%
\begin{Usage}
\begin{verbatim}
make_model_matrices(parsed_formula, data)
\end{verbatim}
\end{Usage}
%
\begin{Arguments}
\begin{ldescription}
\item[\code{parsed\_formula}] Parsed formula object from parse\_formula()

\item[\code{data}] Data frame
\end{ldescription}
\end{Arguments}
%
\begin{Value}
List with X (fixed effects), Z (random effects), and grouping info
\end{Value}
\HeaderA{marginal\_utils}{Utility Functions for Marginal Predictions}{marginal.Rul.utils}
\keyword{internal}{marginal\_utils}
%
\begin{Description}
Internal functions for computing population-averaged (marginal) predictions
via Monte Carlo integration over random effects distribution
\end{Description}
\HeaderA{mc\_integrate\_fixed\_samples}{Monte Carlo integration for marginal predictions}{mc.Rul.integrate.Rul.fixed.Rul.samples}
\keyword{internal}{mc\_integrate\_fixed\_samples}
%
\begin{Description}
Computes E[g\textasciicircum{}(-1)(X'beta + Z'u)] by averaging over draws from u \textasciitilde{} N(0, Sigma\_u)
\end{Description}
%
\begin{Usage}
\begin{verbatim}
mc_integrate_fixed_samples(X, Z, beta, u_samples, inv_link_fn)
\end{verbatim}
\end{Usage}
%
\begin{Arguments}
\begin{ldescription}
\item[\code{X}] Fixed effects design matrix (n x p)

\item[\code{Z}] Random effects design matrix (n x q) for a single observation

\item[\code{beta}] Fixed effects coefficients (p x 1)

\item[\code{u\_samples}] Matrix of random effects samples (n\_sim x q)

\item[\code{inv\_link\_fn}] Inverse link function
\end{ldescription}
\end{Arguments}
%
\begin{Value}
Vector of marginal predictions (length n)
\end{Value}
\HeaderA{mc\_integrate\_marginal}{Monte Carlo integration for marginal predictions (one sample at a time)}{mc.Rul.integrate.Rul.marginal}
\keyword{internal}{mc\_integrate\_marginal}
%
\begin{Description}
More memory-efficient version that processes samples sequentially
\end{Description}
%
\begin{Usage}
\begin{verbatim}
mc_integrate_marginal(X, Z, beta, Sigma_u, inv_link_fn, n_sim = 1000)
\end{verbatim}
\end{Usage}
%
\begin{Arguments}
\begin{ldescription}
\item[\code{X}] Fixed effects design matrix

\item[\code{Z}] Random effects design matrix

\item[\code{beta}] Fixed effects coefficients

\item[\code{Sigma\_u}] Random effects variance-covariance matrix

\item[\code{inv\_link\_fn}] Inverse link function

\item[\code{n\_sim}] Number of Monte Carlo samples
\end{ldescription}
\end{Arguments}
%
\begin{Value}
List with fit and se
\end{Value}
\HeaderA{multinomial}{Multinomial Family for Unordered Categorical Outcomes}{multinomial}
%
\begin{Description}
Create a family object for baseline-category multinomial logit models
through the unified \code{gllamm()} interface.
\end{Description}
%
\begin{Usage}
\begin{verbatim}
multinomial(reference = NULL)
\end{verbatim}
\end{Usage}
%
\begin{Arguments}
\begin{ldescription}
\item[\code{reference}] Reference category (default: first level)
\end{ldescription}
\end{Arguments}
%
\begin{Value}
A family object of class \code{multinomial\_family}
\end{Value}
%
\begin{Examples}
\begin{ExampleCode}
## Not run: 
fit <- gllamm(choice ~ x + (1 | region), data = d, family = multinomial())

## End(Not run)

\end{ExampleCode}
\end{Examples}
\HeaderA{ordinal}{Ordinal Family for Proportional and Non-Proportional Odds Models}{ordinal}
%
\begin{Description}
Create a family object for ordinal regression models with various link functions
\end{Description}
%
\begin{Usage}
\begin{verbatim}
ordinal(
  link = c("logit", "probit", "acl", "crl_forward", "crl_backward", "ppo")
)
\end{verbatim}
\end{Usage}
%
\begin{Arguments}
\begin{ldescription}
\item[\code{link}] Link function to use:
\begin{description}

\item["logit"] Proportional odds (cumulative logit) - default
\item["probit"] Cumulative probit
\item["acl"] Adjacent category logit
\item["crl\_forward"] Forward continuation ratio logit
\item["crl\_backward"] Backward continuation ratio logit
\item["ppo"] Partial proportional odds (non-proportional)

\end{description}

\end{ldescription}
\end{Arguments}
%
\begin{Details}
The ordinal family supports several models for ordered categorical responses:

\strong{Proportional Odds (logit):}
\deqn{P(Y \le k | x) = \frac{1}{1 + \exp(-(\tau_k - x'\beta))}}{}

\strong{Cumulative Probit:}
\deqn{P(Y \le k | x) = \Phi(\tau_k - x'\beta)}{}

\strong{Adjacent Category Logit (ACL):}
Models the log-odds of adjacent categories:
\deqn{\log\frac{P(Y=k)}{P(Y=k-1)} = \alpha_k + x'\beta}{}

\strong{Continuation Ratio Logit (CRL):}
Forward version models sequential decisions:
\deqn{\log\frac{P(Y=k | Y \ge k)}{P(Y>k | Y \ge k)} = \tau_k - x'\beta}{}

Backward version reverses the conditioning.

\strong{Partial Proportional Odds (PPO):}
Relaxes the proportional odds assumption by allowing different
covariate effects per threshold:
\deqn{P(Y \le k | x) = F(\tau_k - x'\beta_k)}{}
\end{Details}
%
\begin{Value}
A family object of class \code{ordinal\_family} with components:
\begin{ldescription}
\item[\code{family}] Character: "ordinal"
\item[\code{link}] Character: name of link function
\item[\code{link\_code}] Integer: numeric code for TMB (1-6)
\end{ldescription}
\end{Value}
%
\begin{Examples}
\begin{ExampleCode}
## Not run: 
# Proportional odds model (default)
family1 <- ordinal()
family2 <- ordinal(link = "logit")

# Adjacent category logit
family3 <- ordinal(link = "acl")

# Partial proportional odds
family4 <- ordinal(link = "ppo")

# Use with gllamm() - recommended interface
fit <- gllamm(rating ~ temp + (1 | judge),
              data = wine,
              family = ordinal(link = "logit"))

# Or use fit_ordinal() directly
fit2 <- fit_ordinal(rating ~ temp + (1 | judge),
                    data = wine,
                    link = "acl")

## End(Not run)

\end{ExampleCode}
\end{Examples}
\HeaderA{parse\_formula}{Parse GLLAMM formula}{parse.Rul.formula}
\keyword{internal}{parse\_formula}
%
\begin{Description}
Parses a formula into fixed effects and random effects components.
Supports lme4-style syntax without requiring lme4.
\end{Description}
%
\begin{Usage}
\begin{verbatim}
parse_formula(formula, data)
\end{verbatim}
\end{Usage}
%
\begin{Arguments}
\begin{ldescription}
\item[\code{formula}] A formula object with syntax: y \textasciitilde{} x + (terms | group)

\item[\code{data}] A data frame containing the variables
\end{ldescription}
\end{Arguments}
%
\begin{Value}
A list with components:
\begin{ldescription}
\item[\code{fixed\_formula}] Formula for fixed effects
\item[\code{random\_terms}] List of random effects specifications
\item[\code{response\_name}] Name of response variable
\end{ldescription}
\end{Value}
\HeaderA{parse\_level\_weights}{Parse observation-level and group-level survey weights}{parse.Rul.level.Rul.weights}
\keyword{internal}{parse\_level\_weights}
%
\begin{Description}
GLLAMM-style level-specific weights: \code{weights} may be a numeric
vector (level-1 / observation weights, as before) or a list with elements
\code{level1} (length n\_obs) and/or \code{level2} (one weight per group,
or per observation but constant within group). Level-2 weights scale each
group's full likelihood contribution including its random-effects prior
(pseudo-likelihood for two-stage sampling designs).
\end{Description}
%
\begin{Usage}
\begin{verbatim}
parse_level_weights(weights, n_obs, groups, n_groups)
\end{verbatim}
\end{Usage}
%
\begin{Arguments}
\begin{ldescription}
\item[\code{weights}] NULL, numeric vector, or list(level1=, level2=)

\item[\code{n\_obs}] Number of observations

\item[\code{groups}] 0-indexed group index per observation

\item[\code{n\_groups}] Number of groups
\end{ldescription}
\end{Arguments}
%
\begin{Value}
list(level1 = numeric(n\_obs), level2 = numeric(n\_groups))
\end{Value}
\HeaderA{parse\_random\_formula}{Parse Random Effects Formula}{parse.Rul.random.Rul.formula}
\keyword{internal}{parse\_random\_formula}
%
\begin{Description}
Internal function to parse lme4-style random effects formulas
\end{Description}
%
\begin{Usage}
\begin{verbatim}
parse_random_formula(formula, data)
\end{verbatim}
\end{Usage}
%
\begin{Arguments}
\begin{ldescription}
\item[\code{formula}] Random effects formula (e.g., \textasciitilde{} (1 | class) + (1 | school))

\item[\code{data}] Data frame containing grouping variables
\end{ldescription}
\end{Arguments}
%
\begin{Value}
List of parsed random effects terms
\end{Value}
\HeaderA{parse\_random\_term}{Parse a single random effects term}{parse.Rul.random.Rul.term}
\keyword{internal}{parse\_random\_term}
%
\begin{Description}
Parse a single random effects term
\end{Description}
%
\begin{Usage}
\begin{verbatim}
parse_random_term(term, data)
\end{verbatim}
\end{Usage}
%
\begin{Arguments}
\begin{ldescription}
\item[\code{term}] Character string with format "(terms | group)"

\item[\code{data}] Data frame for checking variables
\end{ldescription}
\end{Arguments}
%
\begin{Value}
List with parsed random effects structure
\end{Value}
\HeaderA{parse\_single\_random\_term}{Parse a single random effects term}{parse.Rul.single.Rul.random.Rul.term}
\keyword{internal}{parse\_single\_random\_term}
%
\begin{Description}
Parse a single random effects term
\end{Description}
%
\begin{Usage}
\begin{verbatim}
parse_single_random_term(term, data)
\end{verbatim}
\end{Usage}
\HeaderA{plot.gllamm}{Plot diagnostics for GLLAMM models}{plot.gllamm}
%
\begin{Description}
Plot diagnostics for GLLAMM models
\end{Description}
%
\begin{Usage}
\begin{verbatim}
## S3 method for class 'gllamm'
plot(x, which = c(1, 2, 3, 5), ...)
\end{verbatim}
\end{Usage}
%
\begin{Arguments}
\begin{ldescription}
\item[\code{x}] A gllamm object

\item[\code{which}] Which plots to produce (1-6)

\item[\code{...}] Additional arguments
\end{ldescription}
\end{Arguments}
\HeaderA{plot.gllamm\_irt}{Plot IRT Model Diagnostics}{plot.gllamm.Rul.irt}
%
\begin{Description}
Create diagnostic plots for IRT models including item characteristic curves,
item information functions, test information, and ability distributions
\end{Description}
%
\begin{Usage}
\begin{verbatim}
## S3 method for class 'gllamm_irt'
plot(x, which = 1:4, items = NULL, ...)
\end{verbatim}
\end{Usage}
%
\begin{Arguments}
\begin{ldescription}
\item[\code{x}] A gllamm\_irt object

\item[\code{which}] Which plots to produce (1=ICC, 2=IIF, 3=TIF, 4=Ability)

\item[\code{items}] Which items to plot (default: first 6 items)

\item[\code{...}] Additional arguments passed to plotting functions
\end{ldescription}
\end{Arguments}
%
\begin{Details}
Plot types:
\begin{itemize}

\item{} \code{which = 1}: Item Characteristic Curves (ICC)
\item{} \code{which = 2}: Item Information Functions (IIF)
\item{} \code{which = 3}: Test Information Function (TIF)
\item{} \code{which = 4}: Person Ability Distribution

\end{itemize}

\end{Details}
%
\begin{Examples}
\begin{ExampleCode}
## Not run: 
# Fit 2PL model
fit <- fit_irt(responses, model = "2PL")

# Plot all diagnostics for items 1-3
plot(fit, which = 1:4, items = 1:3)

# Plot only ICCs
plot(fit, which = 1, items = 1:5)

## End(Not run)

\end{ExampleCode}
\end{Examples}
\HeaderA{plot.gllamm\_lca}{Plot Latent Class Analysis Results}{plot.gllamm.Rul.lca}
%
\begin{Description}
Create diagnostic plots for latent class models including class profiles,
item probability heatmaps, and classification summaries
\end{Description}
%
\begin{Usage}
\begin{verbatim}
## S3 method for class 'gllamm_lca'
plot(x, which = 1:3, ...)
\end{verbatim}
\end{Usage}
%
\begin{Arguments}
\begin{ldescription}
\item[\code{x}] A gllamm\_lca object

\item[\code{which}] Which plots to produce (1=profiles, 2=heatmap, 3=classification)

\item[\code{...}] Additional arguments passed to plotting functions
\end{ldescription}
\end{Arguments}
%
\begin{Details}
Plot types:
\begin{itemize}

\item{} \code{which = 1}: Class Profiles (line plot of item probabilities by class)
\item{} \code{which = 2}: Item Probability Heatmap
\item{} \code{which = 3}: Classification Summary (barplot of class assignments)

\end{itemize}

\end{Details}
%
\begin{Examples}
\begin{ExampleCode}
## Not run: 
# Fit LCA model
fit <- fit_lca(indicators, nclass = 3)

# Plot all diagnostics
plot(fit, which = 1:3)

# Plot only class profiles
plot(fit, which = 1)

## End(Not run)

\end{ExampleCode}
\end{Examples}
\HeaderA{plot.gllamm\_ordinal}{Plot Ordinal Regression Model Diagnostics}{plot.gllamm.Rul.ordinal}
%
\begin{Description}
Create diagnostic plots for ordinal models including cumulative probabilities,
category probabilities, threshold parameters, and covariate effects
\end{Description}
%
\begin{Usage}
\begin{verbatim}
## S3 method for class 'gllamm_ordinal'
plot(x, which = 1:3, covariate = NULL, covariate_values = NULL, ...)
\end{verbatim}
\end{Usage}
%
\begin{Arguments}
\begin{ldescription}
\item[\code{x}] A gllamm\_ordinal object

\item[\code{which}] Which plots to produce (1=cumulative, 2=category, 3=thresholds, 4=effects)

\item[\code{covariate}] Name of covariate to plot (default: first non-intercept covariate)

\item[\code{covariate\_values}] Optional vector of covariate values to plot (default: -2 to 2)

\item[\code{...}] Additional arguments passed to plotting functions
\end{ldescription}
\end{Arguments}
%
\begin{Details}
Plot types:
\begin{itemize}

\item{} \code{which = 1}: Cumulative Probabilities P(Y <= k) vs covariate
\item{} \code{which = 2}: Category Probabilities P(Y = k) vs covariate
\item{} \code{which = 3}: Threshold Parameters on latent scale
\item{} \code{which = 4}: Covariate Effects (shows non-proportional effects for PPO)

\end{itemize}

\end{Details}
%
\begin{Examples}
\begin{ExampleCode}
## Not run: 
# Fit ordinal model
fit <- fit_ordinal(rating ~ temp + contact + (1 | judge),
                   data = wine, link = "logit")

# Plot all diagnostics for 'temp' covariate
plot(fit, which = 1:4, covariate = "temp")

# Plot only cumulative probabilities
plot(fit, which = 1, covariate = "contact")

## End(Not run)

\end{ExampleCode}
\end{Examples}
\HeaderA{plot\_ability\_distribution\_irt}{Plot Person Ability Distribution}{plot.Rul.ability.Rul.distribution.Rul.irt}
\keyword{internal}{plot\_ability\_distribution\_irt}
%
\begin{Description}
Plot Person Ability Distribution
\end{Description}
%
\begin{Usage}
\begin{verbatim}
plot_ability_distribution_irt(x, ...)
\end{verbatim}
\end{Usage}
\HeaderA{plot\_category\_probs\_ordinal}{Plot Category Probabilities}{plot.Rul.category.Rul.probs.Rul.ordinal}
\keyword{internal}{plot\_category\_probs\_ordinal}
%
\begin{Description}
Plot Category Probabilities
\end{Description}
%
\begin{Usage}
\begin{verbatim}
plot_category_probs_ordinal(x, covariate, covariate_values, ...)
\end{verbatim}
\end{Usage}
\HeaderA{plot\_classification\_lca}{Plot Classification Summary}{plot.Rul.classification.Rul.lca}
\keyword{internal}{plot\_classification\_lca}
%
\begin{Description}
Plot Classification Summary
\end{Description}
%
\begin{Usage}
\begin{verbatim}
plot_classification_lca(x, ...)
\end{verbatim}
\end{Usage}
\HeaderA{plot\_classification\_uncertainty}{Plot Individual Classification Uncertainty}{plot.Rul.classification.Rul.uncertainty}
%
\begin{Description}
Visualize posterior probabilities for individual cases
\end{Description}
%
\begin{Usage}
\begin{verbatim}
plot_classification_uncertainty(
  x,
  cases = 1:min(20, nrow(x$posterior)),
  sort_by = c("entropy", "modal", "index"),
  ...
)
\end{verbatim}
\end{Usage}
%
\begin{Arguments}
\begin{ldescription}
\item[\code{x}] A gllamm\_lca object

\item[\code{cases}] Which cases to plot (default: first 20)

\item[\code{sort\_by}] Sort cases by: "entropy" (default), "modal", or "index"

\item[\code{...}] Additional arguments
\end{ldescription}
\end{Arguments}
%
\begin{Examples}
\begin{ExampleCode}
## Not run: 
fit <- fit_lca(data, nclass = 3)
plot_classification_uncertainty(fit, cases = 1:30)

## End(Not run)

\end{ExampleCode}
\end{Examples}
\HeaderA{plot\_class\_profiles\_lca}{Plot Class Profiles}{plot.Rul.class.Rul.profiles.Rul.lca}
\keyword{internal}{plot\_class\_profiles\_lca}
%
\begin{Description}
Plot Class Profiles
\end{Description}
%
\begin{Usage}
\begin{verbatim}
plot_class_profiles_lca(x, ...)
\end{verbatim}
\end{Usage}
\HeaderA{plot\_covariate\_effects\_ordinal}{Plot Covariate Effects}{plot.Rul.covariate.Rul.effects.Rul.ordinal}
\keyword{internal}{plot\_covariate\_effects\_ordinal}
%
\begin{Description}
Plot Covariate Effects
\end{Description}
%
\begin{Usage}
\begin{verbatim}
plot_covariate_effects_ordinal(x, covariate, ...)
\end{verbatim}
\end{Usage}
\HeaderA{plot\_cumulative\_probs\_ordinal}{Plot Cumulative Probabilities}{plot.Rul.cumulative.Rul.probs.Rul.ordinal}
\keyword{internal}{plot\_cumulative\_probs\_ordinal}
%
\begin{Description}
Plot Cumulative Probabilities
\end{Description}
%
\begin{Usage}
\begin{verbatim}
plot_cumulative_probs_ordinal(x, covariate, covariate_values, ...)
\end{verbatim}
\end{Usage}
\HeaderA{plot\_icc\_irt}{Plot Item Characteristic Curves}{plot.Rul.icc.Rul.irt}
\keyword{internal}{plot\_icc\_irt}
%
\begin{Description}
Plot Item Characteristic Curves
\end{Description}
%
\begin{Usage}
\begin{verbatim}
plot_icc_irt(x, items, ...)
\end{verbatim}
\end{Usage}
\HeaderA{plot\_iif\_irt}{Plot Item Information Functions}{plot.Rul.iif.Rul.irt}
\keyword{internal}{plot\_iif\_irt}
%
\begin{Description}
Plot Item Information Functions
\end{Description}
%
\begin{Usage}
\begin{verbatim}
plot_iif_irt(x, items, ...)
\end{verbatim}
\end{Usage}
\HeaderA{plot\_irt}{Plotting Functions for IRT Models}{plot.Rul.irt}
%
\begin{Description}
Visualize item response theory models
\end{Description}
\HeaderA{plot\_item\_covariates}{Plot Item Covariate Effects}{plot.Rul.item.Rul.covariates}
%
\begin{Description}
Visualize the relationship between item covariates and item parameters
\end{Description}
%
\begin{Usage}
\begin{verbatim}
plot_item_covariates(
  object,
  covariate,
  parameter = c("difficulty", "discrimination"),
  ...
)
\end{verbatim}
\end{Usage}
%
\begin{Arguments}
\begin{ldescription}
\item[\code{object}] A gllamm\_eirt object

\item[\code{covariate}] Name of covariate to plot

\item[\code{parameter}] Which parameter to plot: "difficulty" or "discrimination"

\item[\code{...}] Additional arguments passed to plot
\end{ldescription}
\end{Arguments}
%
\begin{Examples}
\begin{ExampleCode}
## Not run: 
fit <- fit_eirt(responses, item_data,
                difficulty_formula = ~ word_freq)

plot_item_covariates(fit, covariate = "word_freq")

## End(Not run)

\end{ExampleCode}
\end{Examples}
\HeaderA{plot\_item\_probabilities\_lca}{Plot Item Probability Heatmap}{plot.Rul.item.Rul.probabilities.Rul.lca}
\keyword{internal}{plot\_item\_probabilities\_lca}
%
\begin{Description}
Plot Item Probability Heatmap
\end{Description}
%
\begin{Usage}
\begin{verbatim}
plot_item_probabilities_lca(x, ...)
\end{verbatim}
\end{Usage}
\HeaderA{plot\_lca}{Plotting Functions for Latent Class Analysis}{plot.Rul.lca}
%
\begin{Description}
Visualize latent class analysis results
\end{Description}
\HeaderA{plot\_ordinal}{Plotting Functions for Ordinal Regression Models}{plot.Rul.ordinal}
%
\begin{Description}
Visualize ordinal regression model results
\end{Description}
\HeaderA{plot\_ordinal\_effects}{Plot Ordinal Model Effects for Multiple Covariates}{plot.Rul.ordinal.Rul.effects}
%
\begin{Description}
Compare effects across multiple covariates in an ordinal model
\end{Description}
%
\begin{Usage}
\begin{verbatim}
plot_ordinal_effects(
  object,
  covariates = NULL,
  sort_by = c("magnitude", "name", "none"),
  ...
)
\end{verbatim}
\end{Usage}
%
\begin{Arguments}
\begin{ldescription}
\item[\code{object}] A gllamm\_ordinal object

\item[\code{covariates}] Vector of covariate names (default: all non-intercept)

\item[\code{sort\_by}] Sort covariates by: "magnitude", "name", or "none"

\item[\code{...}] Additional arguments
\end{ldescription}
\end{Arguments}
%
\begin{Examples}
\begin{ExampleCode}
## Not run: 
fit <- fit_ordinal(rating ~ temp + contact + (1 | judge), data = wine)
plot_ordinal_effects(fit)

## End(Not run)

\end{ExampleCode}
\end{Examples}
\HeaderA{plot\_thresholds\_ordinal}{Plot Threshold Parameters}{plot.Rul.thresholds.Rul.ordinal}
\keyword{internal}{plot\_thresholds\_ordinal}
%
\begin{Description}
Plot Threshold Parameters
\end{Description}
%
\begin{Usage}
\begin{verbatim}
plot_thresholds_ordinal(x, ...)
\end{verbatim}
\end{Usage}
\HeaderA{plot\_tif\_irt}{Plot Test Information Function}{plot.Rul.tif.Rul.irt}
\keyword{internal}{plot\_tif\_irt}
%
\begin{Description}
Plot Test Information Function
\end{Description}
%
\begin{Usage}
\begin{verbatim}
plot_tif_irt(x, ...)
\end{verbatim}
\end{Usage}
\HeaderA{predict.gllamm}{Predict method for GLLAMM models}{predict.gllamm}
%
\begin{Description}
Obtain predictions from a fitted GLLAMM model
\end{Description}
%
\begin{Usage}
\begin{verbatim}
## S3 method for class 'gllamm'
predict(
  object,
  newdata = NULL,
  type = c("response", "link", "random", "marginal"),
  re.form = NULL,
  n_sim = 1000,
  se.fit = FALSE,
  ...
)
\end{verbatim}
\end{Usage}
%
\begin{Arguments}
\begin{ldescription}
\item[\code{object}] A fitted \code{gllamm} object

\item[\code{newdata}] Optional new data frame for predictions. If omitted, fitted
values from the original data are returned.

\item[\code{type}] Type of prediction:
\begin{description}

\item[response] Predictions on the response scale (default)
\item[link] Predictions on the link scale (linear predictor)
\item[random] Random effects only
\item[marginal] Population-averaged predictions (integrating over random effects)

\end{description}


\item[\code{re.form}] Formula for random effects to include. Use \code{NA} or
\code{\textasciitilde{}0} to exclude all random effects (population-level predictions).
Ignored when \code{type = "marginal"}.

\item[\code{n\_sim}] Number of Monte Carlo samples for marginal predictions (default: 1000).
Only used when \code{type = "marginal"}.

\item[\code{se.fit}] Logical; return standard errors for marginal predictions? (default: FALSE).
Only used when \code{type = "marginal"}.

\item[\code{...}] Additional arguments (currently unused)
\end{ldescription}
\end{Arguments}
%
\begin{Value}
A vector of predictions
\end{Value}
%
\begin{Examples}
\begin{ExampleCode}
## Not run: 
fit <- gllamm(y ~ x + (1 | group), data = mydata)

# Fitted values (default - conditional on random effects)
pred1 <- predict(fit)

# Population-level predictions (fixed effects only, u=0)
pred2 <- predict(fit, re.form = NA)

# Marginal predictions (population-averaged, integrating over u)
pred3 <- predict(fit, type = "marginal")

# Marginal predictions with standard errors
pred4 <- predict(fit, type = "marginal", se.fit = TRUE)

# Marginal predictions for new data
pred5 <- predict(fit, newdata = newdata, type = "marginal")

## End(Not run)

\end{ExampleCode}
\end{Examples}
\HeaderA{predict.gllamm\_eirt}{Predict method for EIRT models}{predict.gllamm.Rul.eirt}
%
\begin{Description}
Obtain predictions from a fitted Explanatory IRT model.
\end{Description}
%
\begin{Usage}
\begin{verbatim}
## S3 method for class 'gllamm_eirt'
predict(
  object,
  newdata = NULL,
  type = c("probability", "ability", "difficulty", "discrimination", "marginal"),
  ability = NULL,
  n_sim = 1000,
  ...
)
\end{verbatim}
\end{Usage}
%
\begin{Arguments}
\begin{ldescription}
\item[\code{object}] A fitted EIRT model (class gllamm\_eirt)

\item[\code{newdata}] Optional data frame with item covariates for new items.
Must include all variables used in the difficulty and discrimination
formulas. For polytomous models, new-item predictions use the predicted
item location with step deviations at their population value of zero
(PCM/GPCM) or the threshold regression (LPCM); GRM threshold spacing is
item-specific and cannot be predicted for new items.

\item[\code{type}] Type of prediction:
\begin{description}

\item[probability] Item response probabilities. Dichotomous: persons x
items matrix of P(Y=1). Polytomous: list of persons x categories
matrices, one per item.
\item[ability] Person ability estimates (theta)
\item[difficulty] Predicted item difficulties
\item[discrimination] Predicted item discriminations
\item[marginal] Marginal response probabilities, integrating ability
over its population distribution. Dichotomous: vector of E[P(Y=1)].
Polytomous: items x categories matrix of E[P(Y=k)].

\end{description}


\item[\code{ability}] Optional vector of ability values. If NULL, uses estimated abilities.

\item[\code{n\_sim}] Number of Monte Carlo samples for marginal predictions (default: 1000)

\item[\code{...}] Additional arguments (currently unused)
\end{ldescription}
\end{Arguments}
%
\begin{Value}
Depends on \code{type}
\end{Value}
\HeaderA{predict.gllamm\_irt}{Predict method for IRT models}{predict.gllamm.Rul.irt}
%
\begin{Description}
Obtain predictions from a fitted IRT model
\end{Description}
%
\begin{Usage}
\begin{verbatim}
## S3 method for class 'gllamm_irt'
predict(
  object,
  newdata = NULL,
  type = c("probability", "ability", "marginal"),
  ability = NULL,
  n_sim = 1000,
  ...
)
\end{verbatim}
\end{Usage}
%
\begin{Arguments}
\begin{ldescription}
\item[\code{object}] A fitted IRT model (class gllamm\_irt)

\item[\code{newdata}] Optional specification of items/persons for predictions.
Can be:
\begin{itemize}

\item{} NULL: Predict for all original items
\item{} Numeric vector: Item indices to predict
\item{} Data frame with 'item' column: Specific items

\end{itemize}


\item[\code{type}] Type of prediction:
\begin{description}

\item[probability] Item response probabilities P(Y=1|θ)
\item[ability] Person ability estimates (θ)
\item[marginal] Marginal item response probabilities E[P(Y=1|θ)]

\end{description}


\item[\code{ability}] Optional vector of ability values for which to compute probabilities.
If NULL and type="probability", uses estimated abilities from fitted model.

\item[\code{n\_sim}] Number of Monte Carlo samples for marginal predictions (default: 1000)

\item[\code{...}] Additional arguments (currently unused)
\end{ldescription}
\end{Arguments}
%
\begin{Value}
Depends on \code{type}:
\begin{itemize}

\item{} probability: Matrix of probabilities (n\_persons × n\_items) or vector if ability specified
\item{} ability: Vector of person abilities
\item{} marginal: Vector of marginal probabilities (one per item)

\end{itemize}

\end{Value}
%
\begin{Examples}
\begin{ExampleCode}
## Not run: 
# Fit 2PL model
responses <- matrix(rbinom(1000, 1, 0.6), 100, 10)
fit <- fit_irt(responses, model = "2PL")

# Person abilities
abilities <- predict(fit, type = "ability")

# Item response probabilities for each person
probs <- predict(fit, type = "probability")

# Marginal item response probabilities (population-level)
marg_probs <- predict(fit, type = "marginal")

## End(Not run)

\end{ExampleCode}
\end{Examples}
\HeaderA{predict.gllamm\_multinomial}{Predict method for multinomial models}{predict.gllamm.Rul.multinomial}
%
\begin{Description}
Obtain predictions from a fitted multinomial logit model with random effects
\end{Description}
%
\begin{Usage}
\begin{verbatim}
## S3 method for class 'gllamm_multinomial'
predict(
  object,
  newdata = NULL,
  type = c("class", "probs", "marginal"),
  n_sim = 1000,
  ...
)
\end{verbatim}
\end{Usage}
%
\begin{Arguments}
\begin{ldescription}
\item[\code{object}] A fitted multinomial model

\item[\code{newdata}] Optional new data frame for predictions

\item[\code{type}] Type of prediction:
\begin{description}

\item[class] Predicted class (modal category)
\item[probs] Conditional probabilities for each category
\item[marginal] Marginal probabilities (population-averaged)

\end{description}


\item[\code{n\_sim}] Number of Monte Carlo samples for marginal predictions (default: 1000)

\item[\code{...}] Additional arguments (currently unused)
\end{ldescription}
\end{Arguments}
%
\begin{Value}
Depends on \code{type}:
\begin{itemize}

\item{} class: Vector of predicted classes
\item{} probs: Matrix of probabilities (n\_obs × n\_categories)
\item{} marginal: Matrix of marginal probabilities (n\_obs × n\_categories)

\end{itemize}

\end{Value}
\HeaderA{predict.gllamm\_ordinal}{Predict method for ordinal models}{predict.gllamm.Rul.ordinal}
%
\begin{Description}
Obtain predictions from a fitted ordinal regression model
\end{Description}
%
\begin{Usage}
\begin{verbatim}
## S3 method for class 'gllamm_ordinal'
predict(
  object,
  newdata = NULL,
  type = c("class", "probs", "cumprobs", "marginal"),
  n_sim = 1000,
  ...
)
\end{verbatim}
\end{Usage}
%
\begin{Arguments}
\begin{ldescription}
\item[\code{object}] A fitted ordinal model (class gllamm with ordinal\_family)

\item[\code{newdata}] Optional new data frame for predictions

\item[\code{type}] Type of prediction:
\begin{description}

\item[class] Predicted class (modal category)
\item[probs] Conditional probabilities for each category
\item[cumprobs] Conditional cumulative probabilities
\item[marginal] Marginal probabilities (population-averaged)

\end{description}


\item[\code{n\_sim}] Number of Monte Carlo samples for marginal predictions (default: 1000)

\item[\code{...}] Additional arguments (currently unused)
\end{ldescription}
\end{Arguments}
%
\begin{Value}
Depends on \code{type}:
\begin{itemize}

\item{} class: Vector of predicted classes
\item{} probs: Matrix of probabilities (n\_obs × n\_categories)
\item{} cumprobs: Matrix of cumulative probabilities
\item{} marginal: Matrix of marginal probabilities (n\_obs × n\_categories)

\end{itemize}

\end{Value}
%
\begin{Examples}
\begin{ExampleCode}
## Not run: 
# Fit ordinal model
fit <- gllamm(rating ~ temp + (1 | judge),
              data = wine,
              family = ordinal(link = "logit"))

# Predicted classes
pred_class <- predict(fit, type = "class")

# Conditional probabilities
pred_probs <- predict(fit, type = "probs")

# Marginal probabilities (population-averaged)
pred_marg <- predict(fit, type = "marginal")

## End(Not run)

\end{ExampleCode}
\end{Examples}
\HeaderA{predict.gllamm\_survival}{Predict method for survival models}{predict.gllamm.Rul.survival}
%
\begin{Description}
Obtain predictions from a fitted survival model with random effects
\end{Description}
%
\begin{Usage}
\begin{verbatim}
## S3 method for class 'gllamm_survival'
predict(
  object,
  newdata = NULL,
  type = c("lp", "risk", "survival", "hazard", "marginal_survival", "marginal_hazard"),
  times = NULL,
  n_sim = 1000,
  ...
)
\end{verbatim}
\end{Usage}
%
\begin{Arguments}
\begin{ldescription}
\item[\code{object}] A fitted survival model

\item[\code{newdata}] Optional new data frame for predictions

\item[\code{type}] Type of prediction:
\begin{description}

\item[lp] Linear predictor (η)
\item[risk] Relative risk exp(η)
\item[survival] Survival probability at specified times
\item[hazard] Hazard at specified times
\item[marginal\_survival] Marginal survival probability (population-averaged)
\item[marginal\_hazard] Marginal hazard (population-averaged)

\end{description}


\item[\code{times}] Times at which to evaluate survival/hazard (required for survival/hazard types)

\item[\code{n\_sim}] Number of Monte Carlo samples for marginal predictions (default: 1000)

\item[\code{...}] Additional arguments (currently unused)
\end{ldescription}
\end{Arguments}
%
\begin{Value}
Depends on \code{type}
\end{Value}
\HeaderA{predict\_difficulty}{Predict Item Difficulties from Covariates}{predict.Rul.difficulty}
%
\begin{Description}
Compute predicted difficulties based on item covariate model
\end{Description}
%
\begin{Usage}
\begin{verbatim}
predict_difficulty(object, newdata = NULL)
\end{verbatim}
\end{Usage}
%
\begin{Arguments}
\begin{ldescription}
\item[\code{object}] A gllamm\_eirt object

\item[\code{newdata}] Optional new item data for predictions
\end{ldescription}
\end{Arguments}
%
\begin{Value}
Vector of predicted difficulties (fixed effects only, no residuals)
\end{Value}
%
\begin{Examples}
\begin{ExampleCode}
## Not run: 
fit <- fit_eirt(responses, item_data,
                difficulty_formula = ~ word_freq)

# Predicted difficulties for fitted data
pred_diff <- predict_difficulty(fit)

# Predicted difficulties for new items
new_items <- data.frame(word_freq = c(-1, 0, 1))
pred_new <- predict_difficulty(fit, newdata = new_items)

## End(Not run)

\end{ExampleCode}
\end{Examples}
\HeaderA{predict\_marginal\_gllamm}{Internal function for marginal predictions}{predict.Rul.marginal.Rul.gllamm}
\keyword{internal}{predict\_marginal\_gllamm}
%
\begin{Description}
Computes population-averaged predictions by integrating over random effects
\end{Description}
%
\begin{Usage}
\begin{verbatim}
predict_marginal_gllamm(object, newdata = NULL, n_sim = 1000, se.fit = FALSE)
\end{verbatim}
\end{Usage}
%
\begin{Arguments}
\begin{ldescription}
\item[\code{object}] Fitted gllamm object

\item[\code{newdata}] New data for predictions (NULL = use original data)

\item[\code{n\_sim}] Number of Monte Carlo samples

\item[\code{se.fit}] Return standard errors?
\end{ldescription}
\end{Arguments}
%
\begin{Value}
Vector of marginal predictions, or list with fit and se.fit
\end{Value}
\HeaderA{predict\_marginal\_irt}{Internal function for marginal IRT predictions}{predict.Rul.marginal.Rul.irt}
\keyword{internal}{predict\_marginal\_irt}
%
\begin{Description}
Computes E[P(Y=1|θ)] where θ \textasciitilde{} N(0, σ²\_θ)
\end{Description}
%
\begin{Usage}
\begin{verbatim}
predict_marginal_irt(object, items, n_sim = 1000)
\end{verbatim}
\end{Usage}
%
\begin{Arguments}
\begin{ldescription}
\item[\code{object}] Fitted IRT model

\item[\code{items}] Item indices to predict

\item[\code{n\_sim}] Number of MC samples
\end{ldescription}
\end{Arguments}
%
\begin{Value}
Vector of marginal probabilities (one per item)
\end{Value}
\HeaderA{predict\_marginal\_ordinal}{Internal function for marginal ordinal predictions}{predict.Rul.marginal.Rul.ordinal}
\keyword{internal}{predict\_marginal\_ordinal}
%
\begin{Description}
Internal function for marginal ordinal predictions
\end{Description}
%
\begin{Usage}
\begin{verbatim}
predict_marginal_ordinal(object, X, Z, n_sim = 1000)
\end{verbatim}
\end{Usage}
%
\begin{Arguments}
\begin{ldescription}
\item[\code{object}] Fitted ordinal model

\item[\code{X}] Fixed effects design matrix

\item[\code{Z}] Random effects design matrix

\item[\code{n\_sim}] Number of MC samples
\end{ldescription}
\end{Arguments}
%
\begin{Value}
Matrix of marginal probabilities (n\_obs × n\_categories)
\end{Value}
\HeaderA{print.binomial\_family}{Print method for binomial family}{print.binomial.Rul.family}
\keyword{internal}{print.binomial\_family}
%
\begin{Description}
Print method for binomial family
\end{Description}
%
\begin{Usage}
\begin{verbatim}
## S3 method for class 'binomial_family'
print(x, ...)
\end{verbatim}
\end{Usage}
\HeaderA{print.dif\_analysis}{Print DIF analysis results}{print.dif.Rul.analysis}
%
\begin{Description}
Print DIF analysis results
\end{Description}
%
\begin{Usage}
\begin{verbatim}
## S3 method for class 'dif_analysis'
print(x, ...)
\end{verbatim}
\end{Usage}
%
\begin{Arguments}
\begin{ldescription}
\item[\code{x}] Object of class dif\_analysis

\item[\code{...}] Additional arguments (not used)
\end{ldescription}
\end{Arguments}
\HeaderA{print.eirt\_comparison}{Print EIRT comparison}{print.eirt.Rul.comparison}
\keyword{internal}{print.eirt\_comparison}
%
\begin{Description}
Print EIRT comparison
\end{Description}
%
\begin{Usage}
\begin{verbatim}
## S3 method for class 'eirt_comparison'
print(x, ...)
\end{verbatim}
\end{Usage}
\HeaderA{print.eirt\_test}{Print EIRT test results}{print.eirt.Rul.test}
\keyword{internal}{print.eirt\_test}
%
\begin{Description}
Print EIRT test results
\end{Description}
%
\begin{Usage}
\begin{verbatim}
## S3 method for class 'eirt_test'
print(x, ...)
\end{verbatim}
\end{Usage}
\HeaderA{print.fit\_statistics}{Print Method for Fit Statistics}{print.fit.Rul.statistics}
%
\begin{Description}
Print Method for Fit Statistics
\end{Description}
%
\begin{Usage}
\begin{verbatim}
## S3 method for class 'fit_statistics'
print(x, ...)
\end{verbatim}
\end{Usage}
%
\begin{Arguments}
\begin{ldescription}
\item[\code{x}] A fit\_statistics object

\item[\code{...}] Additional arguments
\end{ldescription}
\end{Arguments}
\HeaderA{print.gllamm\_eirt}{Print EIRT model results}{print.gllamm.Rul.eirt}
%
\begin{Description}
Print EIRT model results
\end{Description}
%
\begin{Usage}
\begin{verbatim}
## S3 method for class 'gllamm_eirt'
print(x, ...)
\end{verbatim}
\end{Usage}
%
\begin{Arguments}
\begin{ldescription}
\item[\code{x}] Object of class gllamm\_eirt

\item[\code{...}] Additional arguments
\end{ldescription}
\end{Arguments}
\HeaderA{print.ordinal\_family}{Print method for ordinal family}{print.ordinal.Rul.family}
\keyword{internal}{print.ordinal\_family}
%
\begin{Description}
Print method for ordinal family
\end{Description}
%
\begin{Usage}
\begin{verbatim}
## S3 method for class 'ordinal_family'
print(x, ...)
\end{verbatim}
\end{Usage}
\HeaderA{print.po\_test}{Print method for proportional odds test}{print.po.Rul.test}
\keyword{internal}{print.po\_test}
%
\begin{Description}
Print method for proportional odds test
\end{Description}
%
\begin{Usage}
\begin{verbatim}
## S3 method for class 'po_test'
print(x, ...)
\end{verbatim}
\end{Usage}
\HeaderA{ranef}{Generic ranef}{ranef}
%
\begin{Description}
Generic function to extract random effects from fitted models
\end{Description}
%
\begin{Usage}
\begin{verbatim}
ranef(object, ...)

ranef(object, ...)
\end{verbatim}
\end{Usage}
%
\begin{Arguments}
\begin{ldescription}
\item[\code{object}] A fitted model object

\item[\code{...}] Additional arguments passed to methods
\end{ldescription}
\end{Arguments}
%
\begin{Value}
Random effects (method-specific format)
\end{Value}
\HeaderA{ranef.gllamm}{Extract random effects}{ranef.gllamm}
%
\begin{Description}
Extract random effects
\end{Description}
%
\begin{Usage}
\begin{verbatim}
## S3 method for class 'gllamm'
ranef(object, ...)
\end{verbatim}
\end{Usage}
%
\begin{Arguments}
\begin{ldescription}
\item[\code{object}] A gllamm object

\item[\code{...}] Additional arguments (ignored)
\end{ldescription}
\end{Arguments}
%
\begin{Value}
List of random effects by group
\end{Value}
\HeaderA{ranef.gllamm\_irt\_multilevel}{Extract Random Effects from Multi-Level IRT Models}{ranef.gllamm.Rul.irt.Rul.multilevel}
%
\begin{Description}
Extract estimated random effects (group-level deviations) from fitted models
\end{Description}
%
\begin{Usage}
\begin{verbatim}
## S3 method for class 'gllamm_irt_multilevel'
ranef(object, level = NULL, ...)
\end{verbatim}
\end{Usage}
%
\begin{Arguments}
\begin{ldescription}
\item[\code{object}] A fitted multi-level IRT model

\item[\code{level}] Which random effect level to extract. If NULL, returns all levels.

\item[\code{...}] Additional arguments (not used)
\end{ldescription}
\end{Arguments}
%
\begin{Value}
A named vector or matrix of random effects
\end{Value}
%
\begin{Examples}
\begin{ExampleCode}
## Not run: 
# Extract class effects
ranef(fit, level = "class")

# Extract all random effects
ranef(fit)

## End(Not run)

\end{ExampleCode}
\end{Examples}
\HeaderA{rmvnorm\_chol}{Draw samples from a zero-mean multivariate normal}{rmvnorm.Rul.chol}
\keyword{internal}{rmvnorm\_chol}
%
\begin{Description}
Cholesky-based sampler so MASS can remain in Suggests.
\end{Description}
%
\begin{Usage}
\begin{verbatim}
rmvnorm_chol(n, Sigma)
\end{verbatim}
\end{Usage}
%
\begin{Arguments}
\begin{ldescription}
\item[\code{n}] Number of samples

\item[\code{Sigma}] Variance-covariance matrix
\end{ldescription}
\end{Arguments}
%
\begin{Value}
n x ncol(Sigma) matrix of draws
\end{Value}
\HeaderA{sandwich\_vcov\_gllamm}{Cluster-robust (sandwich) covariance for GLLAMM fixed effects}{sandwich.Rul.vcov.Rul.gllamm}
\keyword{internal}{sandwich\_vcov\_gllamm}
%
\begin{Description}
Computes the cluster-robust covariance \eqn{H^{-1} M H^{-1}}{} where
\eqn{H}{} is the observed information of the outer parameters and
\eqn{M = \sum_j s_j s_j'}{} accumulates per-cluster score vectors. Because
clusters are independent, the Laplace marginal log-likelihood decomposes
by cluster; each cluster's score is evaluated by rebuilding its objective
on the cluster's rows and differentiating at the full-data estimates
(with the cluster's random effects re-profiled internally).
\end{Description}
%
\begin{Usage}
\begin{verbatim}
sandwich_vcov_gllamm(object)
\end{verbatim}
\end{Usage}
%
\begin{Arguments}
\begin{ldescription}
\item[\code{object}] A two-level GLMM fitted via \code{gllamm()} (gaussian,
binomial, poisson, or Gamma family, single random-effects term)
\end{ldescription}
\end{Arguments}
%
\begin{Value}
Covariance matrix for the full outer parameter vector, with a
\code{"fixed"} attribute holding the fixed-effects block.
\end{Value}
%
\begin{SeeAlso}
\code{\LinkA{vcov.gllamm}{vcov.gllamm}} with \code{type = "sandwich"}
\end{SeeAlso}
\HeaderA{simulate.gllamm}{Simulate from a GLLAMM model}{simulate.gllamm}
%
\begin{Description}
Simulate response data from a fitted GLLAMM model. Random effects are
drawn fresh from their estimated distribution for every replicate
(population-level simulation), both for the original data and for
\code{newdata}.
\end{Description}
%
\begin{Usage}
\begin{verbatim}
## S3 method for class 'gllamm'
simulate(object, nsim = 1, seed = NULL, newdata = NULL, ...)
\end{verbatim}
\end{Usage}
%
\begin{Arguments}
\begin{ldescription}
\item[\code{object}] A fitted \code{gllamm} object

\item[\code{nsim}] Number of simulations (default: 1)

\item[\code{seed}] Optional random seed (stored in the \code{"seed"} attribute,
following the \code{\LinkA{simulate}{simulate}} contract)

\item[\code{newdata}] Optional new data frame containing the covariates and
grouping variables of the model formula

\item[\code{...}] Additional arguments (currently unused)
\end{ldescription}
\end{Arguments}
%
\begin{Value}
A data frame with \code{nsim} columns, one simulated response
vector per column, with a \code{"seed"} attribute.
\end{Value}
%
\begin{Examples}
\begin{ExampleCode}
## Not run: 
fit <- gllamm(y ~ x + (1 | group), data = mydata)

# Single simulation
sim1 <- simulate(fit)

# Multiple simulations
sim10 <- simulate(fit, nsim = 10)

## End(Not run)

\end{ExampleCode}
\end{Examples}
\HeaderA{simulate.gllamm\_ordinal}{Simulate from a fitted ordinal model}{simulate.gllamm.Rul.ordinal}
%
\begin{Description}
Draws fresh random effects from their estimated distribution for each
replicate (population-level simulation) and samples categories from the
implied response distribution.
\end{Description}
%
\begin{Usage}
\begin{verbatim}
## S3 method for class 'gllamm_ordinal'
simulate(object, nsim = 1, seed = NULL, newdata = NULL, ...)
\end{verbatim}
\end{Usage}
%
\begin{Arguments}
\begin{ldescription}
\item[\code{object}] A fitted \code{gllamm\_ordinal} object

\item[\code{nsim}] Number of simulations (default: 1)

\item[\code{seed}] Optional random seed (stored in the \code{"seed"} attribute)

\item[\code{newdata}] Optional new data frame with the model covariates and
grouping variables

\item[\code{...}] Additional arguments (currently unused)
\end{ldescription}
\end{Arguments}
%
\begin{Value}
A data frame with \code{nsim} columns of simulated category
responses (integer codes 1..K), with a \code{"seed"} attribute.
\end{Value}
\HeaderA{summary.dif\_analysis}{Summary of DIF analysis}{summary.dif.Rul.analysis}
%
\begin{Description}
Summary of DIF analysis
\end{Description}
%
\begin{Usage}
\begin{verbatim}
## S3 method for class 'dif_analysis'
summary(object, ...)
\end{verbatim}
\end{Usage}
%
\begin{Arguments}
\begin{ldescription}
\item[\code{object}] Object of class dif\_analysis

\item[\code{...}] Additional arguments (not used)
\end{ldescription}
\end{Arguments}
\HeaderA{summary.gllamm\_eirt}{Summary of EIRT model}{summary.gllamm.Rul.eirt}
%
\begin{Description}
Summary of EIRT model
\end{Description}
%
\begin{Usage}
\begin{verbatim}
## S3 method for class 'gllamm_eirt'
summary(object, ...)
\end{verbatim}
\end{Usage}
%
\begin{Arguments}
\begin{ldescription}
\item[\code{object}] Object of class gllamm\_eirt

\item[\code{...}] Additional arguments
\end{ldescription}
\end{Arguments}
\HeaderA{test\_item\_covariates}{Test Item Covariate Effects}{test.Rul.item.Rul.covariates}
%
\begin{Description}
Test whether adding item covariates significantly improves model fit
\end{Description}
%
\begin{Usage}
\begin{verbatim}
test_item_covariates(
  response_matrix,
  item_data,
  difficulty_formula = ~1,
  discrimination_formula = ~1,
  model = c("Rasch", "2PL", "GRM"),
  ...
)
\end{verbatim}
\end{Usage}
%
\begin{Arguments}
\begin{ldescription}
\item[\code{response\_matrix}] Matrix of item responses

\item[\code{item\_data}] Data frame of item-level covariates

\item[\code{difficulty\_formula}] Formula for difficulty regression

\item[\code{discrimination\_formula}] Formula for discrimination regression

\item[\code{model}] IRT model type

\item[\code{...}] Additional arguments passed to fit\_eirt
\end{ldescription}
\end{Arguments}
%
\begin{Value}
A list with:
\begin{ldescription}
\item[\code{full\_model}] EIRT model with covariates
\item[\code{null\_model}] EIRT model without covariates (intercept only)
\item[\code{comparison}] Model comparison results
\end{ldescription}
\end{Value}
%
\begin{Examples}
\begin{ExampleCode}
## Not run: 
item_data <- data.frame(
  word_freq = rnorm(20),
  length = rpois(20, 5)
)

# Test if word frequency matters
result <- test_item_covariates(
  responses,
  item_data,
  difficulty_formula = ~ word_freq + length,
  model = "Rasch"
)

print(result$comparison)

## End(Not run)

\end{ExampleCode}
\end{Examples}
\HeaderA{test\_proportional\_odds}{Test Proportional Odds Assumption}{test.Rul.proportional.Rul.odds}
%
\begin{Description}
Perform a likelihood ratio test of the proportional odds assumption
by comparing a proportional odds model to a partial proportional odds model
\end{Description}
%
\begin{Usage}
\begin{verbatim}
test_proportional_odds(object, data = NULL)
\end{verbatim}
\end{Usage}
%
\begin{Arguments}
\begin{ldescription}
\item[\code{object}] A fitted ordinal regression model (gllamm\_ordinal)

\item[\code{data}] The data frame used to fit the model (required when the
fitted object does not store its data)
\end{ldescription}
\end{Arguments}
%
\begin{Details}
The proportional odds assumption states that the effect of covariates
is the same across all thresholds. This function fits a partial proportional
odds (PPO) model where each threshold can have different covariate effects
and tests whether this provides a significantly better fit.

The test statistic is:
\deqn{LRT = 2(logLik_{PPO} - logLik_{PO})}{}

which follows a chi-squared distribution with degrees of freedom equal to
the difference in number of parameters.

If p < 0.05, the proportional odds assumption is rejected, suggesting
that covariate effects vary across thresholds.
\end{Details}
%
\begin{Value}
An object of class \code{po\_test} with components:
\begin{ldescription}
\item[\code{statistic}] Likelihood ratio test statistic
\item[\code{df}] Degrees of freedom for the test
\item[\code{p\_value}] P-value from chi-squared distribution
\item[\code{conclusion}] Text interpretation of the test result
\item[\code{models}] List containing the base and PPO models
\end{ldescription}
\end{Value}
%
\begin{Note}
This function currently only works for models with logit or probit links.
It will not work with ACL, CRL, or already-PPO models.
\end{Note}
%
\begin{Examples}
\begin{ExampleCode}
## Not run: 
# Fit proportional odds model
fit_po <- fit_ordinal(rating ~ temp + (1 | judge),
                      data = wine, link = "logit")

# Test proportional odds assumption
po_test <- test_proportional_odds(fit_po)
print(po_test)

## End(Not run)

\end{ExampleCode}
\end{Examples}
\HeaderA{validate\_formula}{Validate formula}{validate.Rul.formula}
\keyword{internal}{validate\_formula}
%
\begin{Description}
Check that formula is properly specified
\end{Description}
%
\begin{Usage}
\begin{verbatim}
validate_formula(formula, data)
\end{verbatim}
\end{Usage}
%
\begin{Arguments}
\begin{ldescription}
\item[\code{formula}] Formula object

\item[\code{data}] Data frame
\end{ldescription}
\end{Arguments}
%
\begin{Value}
TRUE if valid, stops with error otherwise
\end{Value}
\HeaderA{validate\_poly\_responses}{Validate and auto-recode polytomous response matrices (shared by the Laplace and EM estimation paths)}{validate.Rul.poly.Rul.responses}
\keyword{internal}{validate\_poly\_responses}
%
\begin{Description}
Validate and auto-recode polytomous response matrices (shared by the
Laplace and EM estimation paths)
\end{Description}
%
\begin{Usage}
\begin{verbatim}
validate_poly_responses(response_matrix)
\end{verbatim}
\end{Usage}
%
\begin{Value}
list(response\_matrix, n\_categories\_per\_item, max\_categories)
\end{Value}
\HeaderA{VarCorr}{Extract Variance Components from Multi-Level Models}{VarCorr}
%
\begin{Description}
Generic function to extract variance components from fitted multi-level models
\end{Description}
%
\begin{Usage}
\begin{verbatim}
VarCorr(x, ...)
\end{verbatim}
\end{Usage}
%
\begin{Arguments}
\begin{ldescription}
\item[\code{x}] A fitted model object

\item[\code{...}] Additional arguments passed to methods
\end{ldescription}
\end{Arguments}
%
\begin{Value}
Variance components (method-specific format)
\end{Value}
\HeaderA{VarCorr.gllamm}{Extract variance components}{VarCorr.gllamm}
%
\begin{Description}
Extract variance components
\end{Description}
%
\begin{Usage}
\begin{verbatim}
## S3 method for class 'gllamm'
VarCorr(x, ...)
\end{verbatim}
\end{Usage}
%
\begin{Arguments}
\begin{ldescription}
\item[\code{x}] A gllamm object

\item[\code{...}] Additional arguments (ignored)
\end{ldescription}
\end{Arguments}
%
\begin{Value}
List of variance-covariance matrices for random effects
\end{Value}
\HeaderA{VarCorr.gllamm\_irt\_multilevel}{Extract Variance Components from Multi-Level IRT Models}{VarCorr.gllamm.Rul.irt.Rul.multilevel}
%
\begin{Description}
Extract variance components from fitted multi-level IRT models
\end{Description}
%
\begin{Usage}
\begin{verbatim}
## S3 method for class 'gllamm_irt_multilevel'
VarCorr(x, ...)
\end{verbatim}
\end{Usage}
%
\begin{Arguments}
\begin{ldescription}
\item[\code{x}] A fitted multi-level IRT model (class gllamm\_irt\_multilevel)

\item[\code{...}] Additional arguments (not used)
\end{ldescription}
\end{Arguments}
%
\begin{Value}
A data frame with variance components
\end{Value}
%
\begin{Examples}
\begin{ExampleCode}
## Not run: 
# Fit multi-level model
fit <- fit_irt(responses, model = "2PL",
               person_data = data, random = ~ (1 | class))

# Extract variance components
VarCorr(fit)

## End(Not run)

\end{ExampleCode}
\end{Examples}
\printindex{}
\end{document}
